\documentclass[ASNA]{USG} %twocolumn
\usepackage{anyfontsize} %

\usepackage{colortbl}   % For coloring table lines
\usepackage{xcolor}     % For color definitions

\usepackage{steinmetz}

\usepackage{pbox}
\usepackage{ulem}



\graphicspath{{./images/}}
\specialissue{Morphologic Hard and Soft Tissue Changes following Tooth
Extraction in Different Maxilla and Mandible Anatomic Areas: Current
Knowledge, Prognostic Factors Analysis, and Clinical Management}

\articletype{ORIGINAL ARTICLE}%
%\subarticletype{Particle Technology and Fluidization}

\received{xx January 20xx}
\revised{xx January 20xx}
\accepted{xx January 20xx}
\journal{Concurrency and Computation: Practice and Experience}
\volume{0}
\copyyear{2026}
\startpage{1}
\articledoi{10.1002/0000}

%\raggedbottom

%%
%\renewcommand\ShowFrameLinethickness{0.15pt}
%\renewcommand*\ShowFrameColor{\color{red}}


\begin{document}
\title{Event-Driven Computing Continuum Orchestration Solver (ECCOS): A Discrete-Event Simulation Framework for Optimal Service Placement and ML-Ready Trace Generation}
%\subtitle{Subtitle}
\transtitle{Event-Driven Computing Continuum Orchestration Solver (ECCOS)}
\subtranstitle{trans-subtitle}
\author[1]{Mireia Jaume}[][\linkedin{https://www.linkedin.com/in/mireia-jaume-aaa295157}]%[\facebook{https://www.facebook.com} \linkedin{https://www.linkedin.com/in/mireia-jaume-aaa295157} \twitter{https://www.twitter.com}]
\author[1]{Isaac Lera}[https://orcid.org/0000-0002-2826-5970]
\author[1]{Carlos Guerrero}[https://orcid.org/0000-0003-2969-0597]
%\author[2]{Claudia Bispo}[https://orcid.org/0000-0001-7037-2721]
%\author[6]{Alexis Conway}[https://orcid.org/0000-0001-7037-2721]

%\author[7]{Uttara Chakraborty}[https://orcid.org/0000-0001-7037-2721]

%\author[8]{Kathleen T. Daniels}[https://orcid.org/0000-0001-7037-2721]

%\author[9]{Gelo de la Cruz}[https://orcid.org/0000-0001-7037-2721]

%\author[10]{Laura Ferrer-Font}[https://orcid.org/0000-0001-7037-2721]

%\author[11]{Andrew Filby}[https://orcid.org/0000-0001-7037-2721]

%\author[12]{David M. Gravano}[https://orcid.org/0000-0001-7037-2721]

%\author[13]{Michael D. Gregory}[https://orcid.org/0000-0001-7037-2721]

%\author[14]{Christopher Hall}[https://orcid.org/0000-0001-7037-2721]

%\author[15]{Christian Kukat}[https://orcid.org/0000-0001-7037-2721]

%\author[16]{André Mozes}[https://orcid.org/0000-0001-7037-2721]

%\author[17]{Diana Ordoñez-Rueda}[https://orcid.org/0000-0001-7037-2721]

%\author[18]{Eva Orlowski-Oliver}[https://orcid.org/0000-0001-7037-2721]

%\author[19]{Isabella Pesce}[https://orcid.org/0000-0001-7037-2721]

%\author[20]{Ziv Porat}[https://orcid.org/0000-0001-7037-2721]

%\author[21]{Nicole J. Poulton}[https://orcid.org/0000-0001-7037-2721]

%\author[22]{Kristen M. Reifel}[https://orcid.org/0000-0001-7037-2721]

%\author[23]{Aja M. Rieger}[https://orcid.org/0000-0001-7037-2721]

%\author[24]{Rachael T. C. Sheridan}[https://orcid.org/0000-0001-7037-2721]

%\author[25]{Gert Van Isterdael}[https://orcid.org/0000-0001-7037-2721]

%\author[14]{Rachael V. Walker}[https://orcid.org/0000-0001-7037-2721]

\authormark{JAUME \textsc{et al.}}
\titlemark{Event-Driven Computing Continuum Orchestration Solver (ECCOS): A Discrete-Event Simulator for Optimal Service Placement and ML-Ready Trace Generation}

%\address[1]{\orgdiv{Key Laboratory for Ecological Metallurgy of Multimetallic Mineral of Ministry of Education, }\orgname{Institution Name, }%
%\orgaddress{\state{State Name, }\country{Country Name}}}

\address{\orgdiv{Computer Science Department, }\orgname{University of Balearic Islands, }%
\orgaddress{\state{Palma, }\country{Spain}}}

%\address[3]{\orgdiv{Department Name, }\orgname{Institution Name, }%
%\orgaddress{\state{State Name, }\country{Country Name}}}

%\address[4]{\orgdiv{Department Name, }\orgname{Institution Name, }%
%\orgaddress{\state{State Name, }\country{Country Name}}}

\corres{Carlos Guerrero  (\email{carlos.guerrero@uib.es}) ~|~ Isaac Lera (\email{isaac.lera@uib.es})}

\editor{\textbf{Academic Editor:} TBD ~|~ \textbf{Guest Editor:} TBD}

\presentaddress{This is sample for present address text this is sample for present address text.}

\fundingInfo{Grant PID2024-158637OB-I00, funded by MICIU/AEI/10.13039/501100011033 and by “ERDF A way of making Europe”, “ERDF/EU”.}

\keywords{Computing Continuum | Discrete-Event Simulation | Synthetic Machine Learing Ready Datasets | Service Placement Problem | Resource Optimization}

\transkeywords{Computing Continuum | Discrete-Event Simulation | Synthetic Machine Learing Ready Datasets | Service Placement Problem | Resource Optimization}

\abstract[ABSTRACT]{\sout{This is a generic template designed for use by multiple journals,
which includes several options for customization. Please
refer the author guidelines and author LaTeX manuscript preparation document for the journal to which you
are submitting in order to confirm that your manuscript will
comply with the journal’s requirements. Please replace this
text with your abstract. This is sample abstract text just for the template display purpose.}}

 \transabstract[transABSTRACT]{This is a generic template designed for use by multiple journals,
which includes several options for customization. Please
refer the author guidelines and author LaTeX manuscript preparation document for the journal to which you
are submitting in order to confirm that your manuscript will
comply with the journal’s requirements. Please replace this
text with your abstract. This is sample abstract text just for the template display purpose.}

\abbr{5-FU, 5-fluorouracil; CFD, computational fluid dynamics; CH, channel; EFS, event-free survival; GBM, glioblastoma multiforme; OS, overall survival; PFS, progression-free
survival; SD, standard deviation.}

\contributed{Hao Zhang and Pengyue D. Guo contributed equally to this study.}

\dedicated{Dedicated to Srivivasa Ramanujano on the occasion of his 125th birth anniversary.}

\copyright{This is an open access article under the terms of the \href{Creative Commons Attribution-NonCommercial}{Creative Commons Attribution-NonCommercial} License, which permits use, distribution and reproduction in any medium, provided the
original~work~is~properly cited and is not used for commercial purposes.
\\[5pt]
  ©  2024 The Author(s) \textit{AIChE Journal} published by Wiley Periodicals LLC on behalf of American Institute of Chemical Engineers.}

%\openaccessstatement

\maketitle
%\afterpage{\aftergroup\restoregeometry}

%\onecolumn


\section{Introduction}
\label{sec:introduction}
% Se contextualiza la transición del Cloud al Edge y la necesidad crítica de orquestación dinámica basada en Inteligencia Artificial.
% Se argumentará cómo la escasez de entornos de prueba realistas frena la investigación en este campo.
% Finalmente, se presentarán de forma explícita las contribuciones del artículo (el framework ECCOS y su taxonomía de eventos matemáticamente fundamentada).


The Computing Continuum (Cloud-Fog-Edge) paradigm has emerged as the foundational infrastructure to support the next generation of ultra-low latency and bandwidth-intensive applications, such as distributed Artificial Intelligence (AI), 5G/6G networks, and the Internet of Things (IoT)~\cite{Al-Dulaimy2024The}. Unlike traditional centralized data centers, the Continuum is characterized by extreme heterogeneity, resource scarcity at the Far-Edge, and high spatiotemporal volatility in both the underlying infrastructure and user demand~\cite{Sajjad2026Service}. In this highly dynamic environment, traditional heuristic approaches for service placement may be insufficient on their own. Consequently, the scientific community has extended and complemented existing optimization strategies with Machine Learning techniques, particularly Reinforcement Learning (RL), to enable more adaptive, real-time orchestration~\cite{Walia2024AI-Empowered}.

However, the advancement of AI applied to Continuum orchestration is hindered by a critical gap in the State of the Art, the scarcity of dynamic evaluation environments and the lack of large-scale training datasets~\cite{Mihaiu2025Resource}. Real-world production traces remain inaccessible due to corporate privacy policies, and existing simulation tools (such as ns-3, iFogSim, or CloudSim) exhibit fundamental shortcomings~\cite{fi17090382,Qi2025A}. First, traditional network simulators operate at the packet level, rendering them computationally unscalable for generating the months of macroscopic data required to train a Machine Learning model. Second, they lack granular and reproducible configurations, forcing researchers to hardcode scenarios directly into the source code. Third, they present a tight coupling between the simulation engine and the optimization policies, hindering the integration of new algorithms or external agents. Finally, they model system volatility in a binary manner (e.g., assuming either perfect operation or total permanent failure), completely ignoring complex transient phenomena such as partial capacity degradations (brownouts), viral demand spikes (flash crowds), or runtime architectural mutations.

To address these challenges, this paper presents \textsc{ECCOS} (\textit{Event-Driven Computing Continuum Orchestration Solver}), a macroscopic simulation environment based on discrete events. ECCOS is explicitly designed from its inception to abstract low-level network complexity and focus on the agile, stochastic, and mathematically grounded generation of stress scenarios. Its primary objective is to provide the scientific community with a standardized tool for evaluating orchestration algorithms and generating massive, labeled synthetic datasets.

To tackle the aforementioned limitations, the main contributions of this work are summarized as follows:

\begin{itemize}
    \item \textbf{Scalable Dataset Generation for Machine Learning:} ECCOS natively produces output traces rich in state variables and metrics, explicitly oriented toward the creation of massive datasets for the continuous training of Machine Learning algorithms (both supervised and RL).
    \item \textbf{Modular Architecture and Decoupled Solvers:} A strictly decoupled design is proposed, separating logical modules (context creation, event engine, state manager, and optimization engine). ECCOS includes several native solvers (Integer Linear Programming, ILP, Greedy heuristics\ldots), but its agnostic interface allows researchers to easily implement and plug in their own custom solvers without altering the framework's core.
    \item \textbf{Fine-grained User Definition via YAML:} A declarative configuration system using YAML files is introduced, enabling users to define hybrid topologies, probabilistic distributions, and system behaviors with extreme granularity, thereby guaranteeing total experimental reproducibility.
    \item \textbf{Rigorous Stochastic Modeling and Event Composition:} The framework features an extensive catalog of configurable events that model real scenario entropy (such as brownouts, temporal user disconnections, and geographical demand shifts). Furthermore, its architecture allows for the seamless composition of new, complex events from the natively implemented base events.
    \item \textbf{Dynamism in the Application Lifecycle:} ECCOS allows the simulation of ``Day 2'' operations, enabling runtime architectural mutations that dynamically alter both the topology and the resource footprint of instantiated microservice chains.
\end{itemize}


\section{State of the Art}
% Se revisarán simuladores existentes como iFogSim, EdgeCloudSim o PureEdgeSim.
% Se resaltarán sus limitaciones: enfoque estático, topologías poco realistas o incapacidad para simular Service Function Chaining (SFC) y degradaciones temporales.
% Se concluirá la sección justificando por qué ECCOS cubre el hueco de la generación masiva de datasets dinámicos.


Over the last decade, numerous simulation frameworks have been developed to model the complex interactions within the IoT-Fog-Cloud continuum. While these tools have provided foundational insights into resource management, latency optimization, and service placement, they present significant structural and functional limitations when addressing modern, data-driven, and highly dynamic microservice architectures~\cite{patel2022fog}.

The current landscape of fog simulators is dominated by discrete-event simulation models, many of which are extensions of the widely used CloudSim toolkit~\cite{5192685}. Prominent examples include iFogSim~\cite{https://doi.org/10.1002/spe.2509}, EdgeCloudSim~\cite{https://doi.org/10.1002/ett.3493}, and PureEdgeSim~\cite{Mechalikh2020PureEdgeSim}. Other notable independent frameworks include YAFS (Yet Another Fog Simulator)~\cite{8758823}, FogNetSim++~\cite{8502760}, and FogTorch$\Pi$~\cite{fogtorch}. Despite their utility, an in-depth analysis of these state-of-the-art frameworks reveals critical shortcomings across several dimensions, which the proposed ECCOS simulator explicitly addresses.

\textbf{Scalable Dataset Generation for Machine Learning.} Modern resource allocation and orchestration strategies increasingly rely on Machine Learning (ML) and Reinforcement Learning (RL) algorithms. However, traditional simulators are not designed with data generation as a primary objective. Simulators such as EdgeCloudSim and YAFS typically export performance metrics via Comma Separated Values (CSV) files intended for post-simulation statistical analysis~\cite{GILL2021100391}. They fail to natively capture the deep, continuous state-action-reward interactions required for RL training. ECCOS overcomes this limitation by natively producing rich output traces loaded with granular state variables and metrics. It is explicitly oriented toward the automated creation of massive, high-dimensional datasets, directly bridging the gap between network simulation and continuous ML algorithm training.

\textbf{Modular Architecture and Decoupled Solvers.} A major drawback of prevalent tools like iFogSim is their monolithic design, which results in a steep learning curve and complicates the implementation of custom algorithms~\cite{MARKUS2020102042}. Modifying resource management policies in these simulators often requires invasive changes to the core framework code. While tools like FogWorkFlowSim attempt to simplify task scheduling evaluation~\cite{fi12050089}, they still lack a strictly decoupled architecture. ECCOS introduces a highly modular design that strictly separates logical components (context creation, event engine, state manager, and optimization engine). Unlike existing tools, ECCOS acts through an agnostic interface that natively includes several solvers while allowing researchers to seamlessly plug in custom solvers without altering the framework's core logic.

\textbf{Fine-grained User Definition and Reproducibility.} The configuration of physical topologies, user behaviors, and application definitions in current simulators is often rigid and fragmented. Simulators like DISSECT-CF-Fog and PureEdgeSim rely on verbose XML files while YAFS and iFogSim utilize JSON or require hardcoded Java/Python programmatic setup. These approaches lack the declarative simplicity and precision needed for complex scenario modeling. ECCOS pioneers a declarative configuration system using YAML files, enabling researchers to define hybrid topologies, probabilistic distributions, and fine-grained system behaviors. This level of extreme granularity not only simplifies the setup of complex scenarios but also guarantees total experimental reproducibility.

\textbf{Rigorous Stochastic Modeling and Event Composition.} Real-world fog environments are characterized by high entropy and volatility. Existing simulators offer limited stochastic modeling; for instance, YAFS supports network failures and node disconnections and FogTorch$\Pi$ uses Monte Carlo methods to simulate deviations in Quality of Service (QoS) profiles~\cite{fi11030055}. However, these frameworks do not allow for the modeling of compounding, complex network phenomena. ECCOS addresses this gap by featuring an extensive catalog of configurable base events that accurately model real network entropy, including temporal user disconnections, localized brownouts, and geographical demand shifts. Furthermore, its architecture supports the seamless composition of new, complex events from these native bases, providing a much more rigorous and realistic stochastic modeling environment than currently available.

\textbf{Dynamism in the Application Lifecycle ("Day 2" Operations).} A critical shortcoming in the current fog simulation literature is the static treatment of application lifecycles. Most simulators, including iFogSim and FogTorch, treat application deployment as a static placement problem. Advanced tools like MobFogSim~\cite{PULIAFITO2020102062} introduce mobility and basic Virtual Machine (VM)/container migration, but they fundamentally maintain the original architecture of the deployed service~\cite{gupta2024simulators}. They lack the capability to model runtime elasticity and architectural evolution. ECCOS distinguishes itself by simulating true "Day 2" operations. It enables runtime architectural mutations, allowing the simulator to dynamically alter both the network topology and the resource footprint of instantiated microservice chains, accurately reflecting the lifecycle dynamism of modern cloud-native applications operating at the edge.


\section{ECCOS Simulator Architecture and Modular Design}
\label{sec:eccos_architecture}
% Visión general del simulador, explicando su naturaleza basada en eventos discretos y su enfoque orientado a la generación de datasets.
% Se introducirá un diagrama de arquitectura general que conecte los cuatro módulos principales descritos a continuación.


We present \textbf{ECCOS} (\textit{Event-Driven Computing Continuum Orchestration Solver}) to overcome the architectural rigidity of traditional edge simulation tools, which frequently interleave stochastic event generation, state accounting, and placement optimization into monolithic execution loops. As illustrated in Figure~\ref{fig:eccos_architecture}, ECCOS employs a modular, highly decoupled design pattern structured around four foundational components: \textsc{Ecco-Architect}, \textsc{Ecco-Generator}, \textsc{Ecco-State}, \textsc{Ecco-Solver}, and \textsc{Ecco-Bridge}. This separation of concerns establishes strict data interfaces and transactional protocols between functional phases, ensuring reproducibility, extensible resource accounting, and pluggable optimization.

\begin{figure}[!t]
\centerline{\includegraphics[width=0.98\linewidth]{images/guerrero_figures/eccos_architecture_fixed.pdf}}
\caption{Modular architecture and transactional data flows of the ECCOS simulation framework, highlighting the decoupled interaction between declarative initialization, discrete-event scheduling, state management, and placement optimization.}
 \label{fig:eccos_architecture}
\end{figure}


The lifecycle of an ECCOS simulation is governed by structured data exchanges and distinct operational responsibilities across its core conceptual modules:
\begin{itemize}
    \item \textbf{\textsc{Ecco-Architect} (Context Configuration Engine):}
    Operating during the \emph{Initialization Phase}, \textsc{Ecco-Architect} ingests declarative user manifests written in YAML (\texttt{Config}). It parses topological definitions, user mobility/disconnection models, and application topologies and profiles to instantiate the foundational data structures: the multi-attribute infrastructure graph (\texttt{InfrastructureSet}), the user population (\texttt{UserSet}), and the microservice DAGs (\texttt{AppSet}). Once validated, \textsc{Ecco-Architect} seeds the initial event chronology into \textsc{Ecco-Generator} and transfers the instantiated structural references to \textsc{Ecco-State}.
    \item \textbf{\textsc{Ecco-Generator} (Discrete-Event Stochastic Engine):}
    Acting as the chronological scheduler of the \emph{Discrete-Time Simulation Loop}, \textsc{Ecco-Generator} maintains the priority queue of stochastic simulation events (e.g., user spatial mobility, transient disconnections, popularity surges, and infrastructure node/edge failures). At each discrete temporal advancement $t = t + \Delta t$, it pops the earliest scheduled event payload and delivers it directly to \textsc{Ecco-State} for structural execution.
    \item \textbf{\textsc{Ecco-State} (System State Manager):}
    \textsc{Ecco-State} serves as the authoritative \emph{Single Source of Truth} and structural execution handler. Upon receiving an event from \textsc{Ecco-Generator}, it mutates the underlying data structures accordingly, updating user coordinates, toggling edge/node operational flags, and adjusting microservice demand profiles. It continuously enforces schema-agnostic resource accounting across static residency attributes and dynamic load attributes. Guided by user-defined trigger policies (\texttt{trigger\_policy}), \textsc{Ecco-State} decides whether to absorb event mutations internally or to fire an optimization request. When triggered, it delivers an immutable, consolidated system state snapshot $\mathcal{S}^{(t)} = (\mathcal{G}^{(t)},\mathcal{A}^{(t)},\mathcal{U}^{(t)})$ to \textsc{Ecco-Solver}.
    \item \textbf{\textsc{Ecco-Solver} (Decoupled Optimization Engine):}
    Operating as a stateless decision engine, \textsc{Ecco-Solver} receives the canonical snapshot $\mathcal{S}^{(t)}$ and formulates the optimal placement problem. Regardless of the underlying algorithmic strategy, the solver computes an optimal mapping of microservices onto physical nodes. It returns a structured placement dictionary back to \textsc{Ecco-State}, which transactionally applies the decision to the persistent infrastructure graph.
    \item \textbf{\textsc{Ecco-Bridge} (Data Integration \& ML Interface):}
    Operating across the \emph{Output and Integration Phase}, \textsc{Ecco-Bridge} acts as the unified telemetry exporter and artificial intelligence interoperability layer. Throughout the simulation loop, it asynchronously consumes two complementary data streams: raw infrastructure telemetry and state observations streamed from \textsc{Ecco-State}, alongside placement decisions, execution latencies, and reward/cost metrics emitted by \textsc{Ecco-Solver}. \textsc{Ecco-Bridge} decouples internal simulation structures from external analytics by supporting both offline dataset serialization and online agent interoperability. In the offline mode, it compiles structured, time-stamped execution traces into standardized, ML-ready formats, facilitating benchmarking, offline regression, and offline Reinforcement Learning training. \sout{In the online mode, it wraps the discrete-time simulation loop within a standardized OpenAI Gymnasium environment interface, exposing the standard \texttt{step(action)} and \texttt{reset()} methods. This enables external Reinforcement Learning agents to query state observations ($s_t$), execute placement actions ($a_t$), and receive instantaneous reward signals ($r_t$) directly from the simulator.}

\end{itemize}

The simulation framework was fully implemented in Python (version 3.13.2), leveraging its robust ecosystem for scientific computing and system modeling. The core logic and internal data structures were developed using standard built-in libraries, whereas intensive data manipulation and numerical computations were optimized through the integration of NumPy (v2.4.4) and Pandas (v3.0.3). To accurately model and analyze the underlying network topologies and complex entity interactions within the simulated infrastructure, the NetworkX library (v3.6.1) was extensively employed. Furthermore, algorithmic decision-making processes, such as resource allocation and service placement tasks, were mathematically formulated and solved utilizing the PuLP optimization framework (v3.3.0). The source code is openly available in a GitHub repository\footnote{\color{red}https://github.com/XXXXXX\color{black}}, which also hosts the result traces generated during the experiments reported in this article (Section~\ref{sect_experimentalevaluationandresults}).

The following subsections describe each of these modules in detail, outlining their internal responsibilities, configuration mechanisms, and role within the overall \textsc{ECCOS} simulation workflow.

\subsection{Ecco-Architect: Context Configuration Engine}
\label{sec:ecco_architect}

The \textsc{Ecco-Architect} module defines the static context from which each \textsc{ECCOS} simulation scenario is instantiated. Through a declarative YAML configuration, users specify the structural and stochastic models that describe the Computing Continuum infrastructure, the spatial distribution and mobility behavior of users, and the composition, resource demands, and popularity of applications~\cite{Akgul2023VFogSim}. This module therefore acts as the scenario-construction layer of the simulator: it generates the initial infrastructure graph, assigns physical and logical attributes to nodes and links, creates the initial user population and mobility profiles, and defines the application workload model used later by the event generator and placement solver.

Many configuration aspects require users to specify statistical distributions together with their corresponding parameters, including resource capacities, mobility properties, workload demands, popularity values, and event inter-arrival times~\cite{Lee2014Spatial, Zanella2023Characterizing, SolisMoreno2014Analysis}. The complete set of distributions currently supported by the simulator, together with their expected YAML parametrization, is provided in the corresponding appendix (Section~\ref{anex:distributions}).

To ensure strict experimental repeatability and prevent stochastic stream interference across architectural layers, the ECCOS framework isolates pseudorandom number generation (PRNG) into dedicated, independently seeded streams for each structural domain: the physical network topology, the application ecosystem, and the user population. Relying on a shared global random state, a procedural adjustments in one component shift the pseudorandom sequence of downstream processes. Thus, each entity set is governed by an autonomous \texttt{NumPy} Generator instance in ECCOS implementation. This domain decoupling supports deterministic repeatability under fixed seeds and execution order across simulation executions and enables combinatorial benchmarking~\cite{Jun2020Vectorization,Casadei2022A}. For instance, researchers can generate diverse user mobility profiles or application demand landscapes over an identical underlying physical network topology without cross-contamination, ensuring that comparative structural evaluations are strictly controlled and reproducible.


\subsubsection{Infrastructure Definition}



Figure~\ref{fig:uml_infradefinition} summarizes the UML model that supports the infrastructure configuration described in this subsection. The discussion below focuses on the main modeling decisions and their role within \textsc{ECCOS}, while the corresponding appendix provides a detailed description of the YAML parameters used to instantiate these infrastructure models (Section~\ref{anex:infrastructuredefinition}).


The infrastructure definition in \textsc{ECCO-Architect} is driven by the statistical and random models used to construct the graph that represents the Computing Continuum. In this graph, devices are modeled as nodes, while network links among them are represented as edges~\cite{Norman2011Connectivity}. Both types of elements are subsequently annotated with configurable dynamic attributes that capture the relevant physical and operational characteristics of real infrastructures, such as computational capacity, memory, latency, bandwidth, or any other resource-related descriptor required by the placement and orchestration logic~\cite{Couto2015Reliability}.

The first step in this process is to specify the number of nodes to be generated and the random or statistical model that determines the structural shape of the graph~\cite{Bu2024Edge}. This choice defines how nodes are interconnected, the resulting topological properties of the simulated infrastructure, and, in combination with the layer-assignment strategy, how each node is mapped onto the cloud, fog, or edge segments of the Computing Continuum. In some cases, this logical hierarchy is produced implicitly during graph construction, whereas in others it is derived afterwards using a model-specific mapping strategy. Consequently, topology generation and layer assignment are complementary mechanisms that jointly determine both the physical connectivity structure and the logical continuum hierarchy~\cite{Armeniakos2025Stochastic}. \textsc{ECCOS} provides five graph-generation models to represent the heterogeneity and geographical distribution of Computing Continuum infrastructures:


\begin{enumerate}
\item \textbf{Scale-Free Topology.}~\cite{Rak2020The} This model is designed to emulate the backbone of the Internet and large-scale telecommunication networks, where a few central nodes act as massive communication hubs. The topology is generated using the Barabási-Albert (BA) model, which relies on a preferential attachment mechanism governed by the number of edges a new node establishes with existing nodes. A centrality-based classification, within user-defined thresholds, are employed to assign nodes to the three layers of the continuum. Nodes exhibiting the highest centrality scores are designated as the Cloud tier, intermediate ones form the Fog, and the vast majority of peripheral nodes with low centrality compose the Edge tier.

\item \textbf{Random Graph (Erd\H{o}s-Rényi) Topology.}~\cite{Bu2024Edge} The simulator implements the Erd\H{o}s-Rényi (ER) probabilistic model to baseline network connectivity and simulate unstructured, highly random ad-hoc meshes infrastructures. In this approach, every potential edge between any two nodes is created independently with a uniform probability. Since this model lacks an inherent hierarchical structure or distinct hubs, the logical mapping to the Computing Continuum also relies on Betweenness Centrality thresholds.

\item \textbf{Spatial (Geometric) Topology.}~\cite{Tosun2023CoRMAC} Aimed at simulating wireless sensor networks, vehicular ad-hoc networks (VANETs), and Far-Edge environments, this model utilizes Random Geometric Graphs. Nodes are uniformly distributed across a two-dimensional Euclidean plane and are interconnected if the distance between them falls below a specified communication radius. The continuum layer assignment also relies on the Betweenness Centrality of the resulting physical graph. Nodes that serendipitously bridge multiple geographical clusters acquire high centrality and are classified as Cloud or Fog gateways, whereas the scattered exterior nodes remain as Edge devices.

\item \textbf{Hierarchical Fat-Tree Topology.}~\cite{Guo2015On} For scenarios focused on modern intra-datacenter networks or highly structured corporate infrastructures, the simulator implements perfectly balanced Tree and highly-redundant Fat-Tree graphs. These topologies are strictly defined by their branching factor and total height. The mapping to the logical continuum is purely deterministic and depth-based (number of hops from the root). The root nodes inherently represents the Cloud datacenter, the intermediate aggregation layers act as Fog switches, and the bottom-tier leaf nodes are strictly designated as the Edge layer.

\item \textbf{Multi-Tier Hybrid Topology.}~\cite{Armeniakos2025Stochastic,Fernando2025On}  To represent the most comprehensive and realistic end-to-end Computing Continuum ecosystem, the simulator introduces a novel composite topology. This model intrinsically builds and stitches three independent subgraphs: (a) Complete Graph to represent the highly interconnected Cloud core; (b) A Barabási-Albert Scale-Free Graph to model the regional Fog gateways; (c) A Random Geometric Graph to simulate the densely populated Edge devices. Once generated, the tiers are geographically synthesized and cross-linked using Euclidean distance minimization, stitching Edge nodes to their closest Fog gateway, and Fog nodes to their nearest Cloud counterpart. In this paradigm, the logical mapping is intrinsic to the generation process, ensuring a structurally robust Cloud-Fog-Edge hierarchy by design.

\end{enumerate}


To realistically characterize the generated nodes and edges, \textsc{ECCOS} relies on configurable attributes that describe the resources and operational properties associated with each infrastructure element. These attributes may be defined for nodes, or edges, and the user can specify as many attributes as required by the target scenario, such as CPU capacity, memory, storage, latency, bandwidth, energy availability, or reliability indicators~\cite{Sefati2025A}. For each attribute, the configuration must define how its initial values are assigned across the generated topology.

Two basic and complementary initialization strategies are supported for those attributes. In the first one, attribute values are sampled independently for each element, regardless of its position in the graph. This mode is appropriate when resources are assumed to be randomly distributed across the infrastructure~\cite{Wang2024Dynamic}. In the second strategy, the assignment is topology-aware and depends on the centrality of each element, so that the structural role of a node or edge influences the resources it receives~\cite{Oldham2018Consistency}. In this case, the centrality value can be used in a direct proportional manner, multiplying the sampled resource values by the centrality score and therefore assigning larger capacities to elements closer to the cloud or to highly central positions. Conversely, an inverse proportional mode multiplies the sampled values by the inverse centrality term, $(1 - \text{centrality})$, producing lower resource values for elements close to the cloud and higher values toward less central regions. The base resource values are generated from an explicit statistical distribution, ensuring that both randomness and topology-dependent structure can be combined in a controlled and reproducible way. In addition, each attribute can define lower and upper clipping bounds to guarantee that, regardless of the selected distribution, the generated values remain within a reasonable and realistic range for that particular resource or property. 


For any of the available topologies, this general resource-assignment behavior can be overridden when a more explicit layer-specific characterization is required. In particular, users may include dedicated configuration blocks for one or more infrastructure layers (cloud, fog, or edge), thereby defining different resource distributions for values initialization and different max and min values for the selected layers~\cite{Kalyani2021A}. In other words, the global attribute-assignment policy, whether homogeneous or centrality-based, acts as a default mechanism, while layer-specific blocks allow the simulator to specialize the resource profile of cloud, fog, or edge nodes whenever the experimental scenario requires differentiated infrastructure capabilities. When explicit configurations are provided for all three layers, the configuration mode can be set to layered, eliminating the need to define a general assignment pattern and allowing each layer configuration to be specified directly.

\begin{figure}[!t]
\centerline{\includegraphics[width=0.98\linewidth]{images/guerrero_figures/eccos_schema_infra_subset.pdf}}
\caption{UML model for infrastructure definition.\label{fig:uml_infradefinition}}
\end{figure}


\subsubsection{User Definition}

The user-configuration schema is summarized in Figure~\ref{fig:uml_userdefinition}, which provides the structural reference for the modeling choices discussed in this subsection. The following paragraphs emphasize the conceptual role of each user-related feature in \textsc{ECCOS}, while the complete YAML-level specification is reported in the corresponding appendix (Section~\ref{anex:userdefinition}).

\begin{figure}[!t]
\centerline{\includegraphics[width=0.98\linewidth]{images/guerrero_figures/eccos_schema_user_subset.pdf}}
\caption{UML model for user definition.\label{fig:uml_userdefinition}}
\end{figure}



To keep the user-application interaction model tractable, \textsc{ECCOS} adopts a simplified initialization logic in which each user requests a single application~\cite{Bettstetter2004Stochastic}. Under this assumption, the config file fixes the exact number of user entities initially present in the simulated system, and the statistical distribution used to instantiate the request rate generated by each user toward its assigned application. This one-user-one-application abstraction avoids unnecessary combinatorial complexity in the initial workload model without reducing the expressiveness of the simulator: more general scenarios in which a user consumes multiple applications can be emulated straightforwardly by generating a larger population of user entities, each representing one user--application demand relationship.



ECCOS supports user mobility during the simulation. To enable this behavior, the configuration must define the physical space in which users move, their initial positions, and the mobility pattern that determines how they change location over time.

The physical space is defined by a two-dimensional plane, with configurable width and height parameters. In addition, the number of equally spaced streets is also user-defined for the specific case of the Manhattan user mobility model. 


For the initial placement of the users in the plane, a first simple spatial distribution model offers a purely strochastic random uniform distribution, which isotropically spawns users across the entire 2D topology space with equal probability~\cite{Han2022Local}. This model serves as a fundamental baseline, emulating widely dispersed and homogeneous user scenarios, such as suburban neighborhoods or rural deployments, where no significant aggregation factors influence network utilization. 

But in real-world urban scenarios, mobile users do not distribute homogeneously across the topology, they tend to aggregate around specific centers of interest such as transit hubs, shopping malls, or stadiums~\cite{Yang2021Comparative}. The simulater implements a second placement model that considers a stochastic Thomas Cluster Process. This spatial model first generates a predefined number of latent "parent points" (or hotspots) across the simulated 2D plane, with their spatial dispersion sampled from a two-dimensional normal distribution. Once these hotspots are established, individual users are assigned to them and their exact coordinates are instantiated using an isotropic two-dimensional Gaussian distribution centered at their respective hotspot~\cite{Saha2016Enriched}. The degree of spatial aggregation is strictly controlled by a variance parameter, which dictates the dispersion of the user population. A tight spread generates highly dense crowds inherently constrained to the hotspot core, whereas a larger spread models a more relaxed distribution, similar to users roaming the surrounding neighborhoods of a central landmark. 


To accurately emulate the inherent heterogeneity of mobile environments, the simulation framework dynamically individualizes the kinetic profile of each user entity during initialization~\cite{Kour2021PERFORMANCE}. The traversal speed and the maximum coverage radius before signal degradation are stochastically sampled from predefined statistical distributions configured in the system~\cite{Bettstetter2001Mobility}. Additionally, the movement patterns are determined by one of the two available mobility models: the Manhattan mobility model strictly constrains user trajectories to an orthogonal grid, mimicking the topology of city streets and avenues~\cite{1208920,Sadiki2025Advanced}; the random direction model provides an unconstrained, isotropic random walk, better suited for simulating open-space environments such as parks, rural areas, or large event venues~\cite{Kerdsri2017Visualization}. In both cases, the probablity on changing the directions of previous movements/crossings is also determined in the YAML config file. 






\subsubsection{Application Definition}

Figure~\ref{fig:uml_applicationdefinition} presents the UML schema used to organize the application-level configuration elements considered in this subsection. Rather than enumerating every YAML field in the main text, the discussion concentrates on how \textsc{ECCOS} models application structure, resource demand, and popularity, with the detailed configuration parameters documented in the corresponding appendix (Section~\ref{anex:applicationdefinition}).


\begin{figure}[!t]
\centerline{\includegraphics[width=0.98\linewidth]{images/guerrero_figures/eccos_schema_app_subset.pdf}}
\caption{UML model for application definition.\label{fig:uml_applicationdefinition}}
\end{figure}



In the context of evaluating service placement strategies in distributed edge environments, the computational workloads, the structural dependencies of applications, and the popularity patterns are influencing features and requires accurately modeling~\cite{Mechalikh2020PureEdgeSim}. 

Regarding the structural dependencies of applications, the simulator considers each application as a set of inter-related services~\cite{Cerny2024From}. The simulator considers a fixed number for the initial number of applications, but the number of services composing each application is sampled from a uniform distribution within a configurable length. The user is able to decide whether applications are handled as a unique block, obtained by encapsulating all services in an indivisible entity (monolithic), or as independent services with different placement decisions (microservices)~\cite{Ali2024EdgeBus}.

In the monolithic case, the simulator treats each application as a single self-contained computational entity whose attributes and features are calculated as the sum of the attributes of its component services. The placement strategy therefore assigns the entire application to a single node. This mode is particularly relevant for simulating legacy systems, lightweight edge scripts, or highly coupled computing tasks~\cite{Oyeniran2024Microservices}.

In the microservice case, applications are treated as distributed sets of loosely coupled, fine-grained services with explicit dataflow dependencies. These dependencies can be configured either as a strict linear sequence, representing Service Function Chaining (SFC)~\cite{Thiruvasagam2021Reliable}, or as directed scale-free graph topologies (DAG)~\cite{Cerny2024From}. In this latter case, the simulator implements a directed scale-free network generation mechanism (typically leveraging preferential attachment algorithms similar to the Barabási–Albert model), dictating the initial network seeds (number of initial nodes) and the degree of connectivity for newly appended nodes. The resulting DAGs reflects the presence of central "hub" services and the intricate branching logic inherent in highly scalable, distributed cloud-native applications~\cite{Zhao2022DAG}.

SFCs are representative of traditional telecommunication network service deployments, such as deep packet inspection or sequential stream processing pipelines, where data flows traverse a predetermined sequence of virtualized network functions~\cite{Adoga2022Network}. On contrary, DAGs model the complex, non-linear dependencies characteristic of modern, web-scale microservice architectures~\cite{Borth2025Directed}.



Regarding computational workloads, applications are formally defined by a set of generalized resource-requirement profiles~\cite{Mechalikh2020PureEdgeSim}, such as IoT telemetry, video streaming\ldots These profiles are randomly selected according to a categorical distribution parameterized by user-defined probabilities. Each profile is characterized by a dynamic set of multidimensional attributes (e.g., CPU, RAM, storage, and network bandwidth). These attributes represent the computational resource demands that services impose on infrastructure nodes. Similarly, communication resource consumption over network edges is modeled for the entire application through a dynamic set of attributes. In both cases, attribute values are initialized according to their underlying stochastic distributions. Attribute names must be mapped to the corresponding infrastructure attributes, whether node- or edge-level, in order to manage constraints such as resource capacity limits.



Applications are also configured according to their popularity, which is determined by a user-defined stochastic distribution. This popularity scheme can be expressed either as a general global request probability or as a geographically localized request probability~\cite{Eramo2017An}. This spatio-temporal popularity model regulates how users interact with services, allowing the framework to emulate realistic traffic-demand distributions for both globally popular and geographically localized applications. Consequently, applications are partitioned into global and local domains through a configurable percentatge of the full number of applications.

The global applications are requested from any region with the same popularity-based probability. On contrary, local applications are assigned with a center of influence represented by a 2D coordinate within the simulation space~\cite{Chen2022Efficient}. The probability of requesting a local application (popularity) is dinamycally adapted for new users created inside the local applicaiton influence radius, increasing the local application popularity according to the distance between the user and the application's center of influence~\cite{Shi2025Dynamic}. In other words, the closer the user is to the center of a local application, the more likely the user is to decide to use it. Conversely, if the user is created beyond the influence radius, this probability boost is not considered, and the application request is governed only by its baseline global popularity as defined by the selected statistical distribution~\cite{Jiang2018User}.


To guarantee that every newly deployed application is immediately subjected to a representative, randomized volume of user requests, the simulator creates a set of new active users for each new application. This allows the service placement algorithms to be evaluated under immediate operational demand.


\subsection{Ecco-Generator: Discrete-Event Stochastic Engine}
\label{sec:ecco_generator}
% Detallará el núcleo estocástico del simulador, responsable de inyectar dinamismo a lo largo del tiempo.
% Se explicarán las mutaciones de aplicaciones, la movilidad tipo Manhattan, y la inyección de caídas y brownouts (modelados con distribuciones Poisson y Weibull).


The \textsc{Ecco-Generator} is the core stochastic engine of the \textsc{ECCOS} framework. Unlike static placement simulators, \textsc{ECCOS} relies on a discrete-event simulation (DES) paradigm to model the extreme volatility and spatiotemporal fluctuations inherent to the Computing Continuum. This module is responsible for injecting entropy into the system across three main entities: Users (Section~\ref{sect_userevents}), Applications (Section~\ref{sect_applicationevents}), and Infrastructure (Section~\ref{sect_infrastructureevents}). 

The simulator operates by generating a list of timestamped events, each one associated with a specific simulation time and with the instances generator (global spawner), or with an already generated instance of a system entity (user, application, node, or edge). These events are then processed in chronological order according to their scheduled simulation time. As a result, the simulation clock does not advance continuously, but discretely: it jumps from one event to the next, updating the system state only when an event is executed~\cite{1611320}. This event-driven mechanism enables ECCOS to efficiently capture dynamic behavior while preserving an explicit temporal ordering of all changes occurring in the simulated Computing Continuum~\cite{MasipBruin2021Managing}.


Each event represents a state transition in the simulated environment and is ultimately reflected in the contextual information managed by \textsc{ECCO-State} (Section~\ref{subsec:ecco_state}). For some event classes, the configuration must specify not only the event identifier and the statistical distribution (Section~\ref{anex:distributions}) that determines when it is triggered (\texttt{frequency}), but also the operational parameters that define how the state transition is applied and its impact (\texttt{impact}) in terms of magnitude, target selection rules, affected resources, recovery behavior, or any additional effect to be required to execute the corresponding transition.

Dynamic simulation events are governed by an independent PRNG stream completely decoupled from the foundational infrastructure, application, and user generation phases (Section~\ref{sec:ecco_architect}). By isolating event-level stochasticity into a dedicated, independently seeded generator, the simulation engine ensures exact execution-to-execution repeatability of dynamic traces across arbitrary time horizons~\cite{Antunes2024Random}. This separation allows researchers to inject distinct stochastic event schedules or operational anomalies into identically initialized infrastructure and workload baselines. Consequently, observed variations in quality-of-service or placement algorithm performance can be unequivocally attributed to the evaluated policies rather than underlying stochastic drift.


In addition to the base events provided by the simulator, \textsc{ECCOS} also supports event encapsulation, allowing users to define composite events that combine multiple elementary state transitions into a higher-level simulation action (Section~\ref{sect_usercompositionofencapsulatedevents}). Importantly, \textsc{ECCOS} does not require all event types to be enabled in a given experiment: events omitted from the YAML configuration are simply disabled. Together, event omission and event encapsulation increase experimental control: researchers can disable irrelevant sources of dynamism by not defining them in the YAML configuration, while also composing higher-level events from existing elementary transitions when more complex behaviors are required.

The general configuration model that captures how both basic simulator events and user-defined encapsulated events are specified is shown in Figure~\ref{fig:uml_events}. This UML view provides a common reference for the following subsections, where the individual event families are described according to their target entity, frequency model, impact parameters, and optional composition mechanisms. Additionally, specific details of the YAML structure for events configuration are documented in several appendices (Sections~\ref{anex_infrastructureevents}, \ref{anex_userevents},  \ref{anex_appsevents},  \ref{anex_globalevents} and  \ref{anex:usercomposedevents}, )

\begin{figure}[!t]
\centerline{\includegraphics[width=0.98\linewidth]{images/guerrero_figures/eccos_schema_global_subset.pdf}}
\caption{UML model for the definition of base events and user-defined encapsulated events.\label{fig:uml_events}}
\end{figure}



\subsubsection{Infrastructure Events}
\label{sect_infrastructureevents}


\begin{center}
\begin{table*}[!ht]%
\caption{Stochastic action catalog and impact configuration schema for the \texttt{infrastructure} domain.\label{tab:infra_actions}}
\begin{tabular*}{\textwidth}{@{\extracolsep\fill}ll l p{5cm}@{\extracolsep\fill}}
\toprule
\multicolumn{2}{@{}l}{\textbf{Action Specification}} & \multicolumn{2}{@{}l}{\textbf{Impact Configuration Schema}} \\\cmidrule{1-2}\cmidrule{3-4}
\textbf{Action Name} & \textbf{Action Description}  & \textbf{Impact Attribute}  & \textbf{Attribute Description \& Semantic Meaning}   \\
\midrule
\multirow{2}{*}{\textbf{\texttt{disable\_node}}} 
 & \multirow{2}{*}{\pbox{4cm}{Completely shuts down a physical server, forcing service re-routing.}} 
 & \texttt{distribution\_to\_enable\_node} & Probability distribution governing downtime duration before auto-revival (\texttt{revive\_node}). \\
\midrule
\multirow{2}{*}{\textbf{\texttt{disable\_edge}}} 
 & \multirow{2}{*}{\pbox{4cm}{Severs a network link, triggering shortest-path recalculations.}} 
 & \texttt{distribution\_to\_enable\_edge} & Probability distribution governing link downtime duration before auto-restoration. \\
\midrule
\multirow{4}{*}{\textbf{\texttt{degrade\_node}}} 
 & \multirow{4}{*}{\pbox{4cm}{Temporarily reduces available computational capacity on a physical server.}} 
 & \texttt{target\_attributes} & List of server resource metrics to degrade (e.g., \texttt{["cpu", "ram"]}). \\
 & & \texttt{p\_loss\_dist} & Probability distribution specifying the percentage reduction of server capacity. \\
 & & \texttt{distribution\_to\_restore\_node} & Distribution governing degradation duration before full capacity is restored. \\
\midrule
\multirow{5}{*}{\textbf{\texttt{congest\_edge}}} 
 & \multirow{5}{*}{\pbox{4cm}{Degrades or inflates network link attributes due to congestion.}} 
 & \texttt{target\_attributes\_loss} & Specific link attributes targeted for capacity loss. \\
  & & \texttt{target\_attributes\_multiplier} & Specific link attributes targeted for delay inflation. \\
 & & \texttt{p\_loss\_dist} & Distribution specifying bandwidth reduction or packet loss percentage. \\
 & & \texttt{multiplier\_dist} & Probability distribution specifying the latency scaling multiplier. \\
 & & \texttt{distribution\_to\_clear\_edge} & Distribution governing congestion duration before normal link metrics return. \\
\bottomrule
\end{tabular*}
\end{table*}
\end{center}


The physical layer of the Computing Continuum is highly volatile. The simulator incorporates events to model availability of the nodes and the edges of the infrastructure (Table~\ref{tab:infra_actions}). \textsc{ECCOS} distinguishes between total failures (blackouts) and temporal degradations (brownouts).


Total failures simulate power outages, hardware crashes, or cable cuts affecting individual infrastructure components~\cite{Pereira2021Availability}. \textsc{ECCOS} models these failures through instance-level events associated with specific entities of type node or edge, rather than as global failure processes over the whole infrastructure class. Accordingly, the edge (\texttt{disable\_edge}) and node (\texttt{disable\_node}) disable events are scheduled independently for each node or edge instance, and their frequency distributions describe the failure behavior of that specific component. When one of these instance-level events is triggered, the affected component is removed from the operational graph, its capabilities become unavailable, and any workloads depending on it must be evacuated or re-placed. The recovery process is also defined at the instance level: once the failure is generated, \textsc{ECCOS} schedules the corresponding restoration event using the impact distribution associated with that event, thereby modeling an individual Mean Time To Repair (MTTR) behavior for the failed node or edge.

In addition, degradation events are also implemented in the simulator. In Edge computing, partial degradations are far more common than total failures. \textsc{ECCOS} models computational throttling (e.g., thermal throttling, OS contention) via degrade node event (\texttt{degrade\_node})~\cite{Bantock2020Mitigating}, and network interference via congest edge event (\texttt{congest\_edge})~\cite{Zhang2022Optimal}. 

During a node degradation event, a node temporarily loses a percentage of its nominal capacity over the user-selected node attributes. Similarly, the network congestion mimicks network propierties. These two events are paired with the recovery events, that restore the initial attributes values and ensure the infrastructure eventually returns to its exact nominal state. This allowes to simulate, for example, situations such as a complex trade-off: whether to migrate away from a congested node at a high cost, or temporarily endure the QoS penalty until the brownout clears.


%\subsubsection{Users, Mobility and Connectivity Dynamics}
\subsubsection{User Events}
\label{sect_userevents}


\begin{center}
\begin{table*}[!ht]%
\caption{Stochastic action catalog and impact configuration schema for the \texttt{user} domain.\label{tab:user_actions}}
\begin{tabular*}{\textwidth}{@{\extracolsep\fill}ll l p{4.8cm}@{\extracolsep\fill}}
\toprule
\multicolumn{2}{@{}l}{\textbf{Action Specification}} & \multicolumn{2}{@{}l}{\textbf{Impact Configuration Schema}} \\\cmidrule{1-2}\cmidrule{3-4}
\textbf{Action Name} & \textbf{Action Description}  & \textbf{Impact Attribute}  & \textbf{Attribute Description \& Semantic Meaning}   \\
\midrule
\textbf{\texttt{remove\_user}} 
 & Permanently terminates a user session. 
 & -- & -- \\
\midrule
\textbf{\texttt{move\_user}} 
 & Displaces user coordinates $(X, Y)$. 
 & -- & -- \\
\midrule
\multirow{2}{*}{\textbf{\texttt{suspend\_user}}} 
 & \multirow{2}{*}{\pbox{4.5cm}{Temporarily halts user session and stops workload generation.}} 
 & \texttt{distribution\_to\_resume\_user} & Distribution governing suspension duration before auto-resumption. \\
\midrule
\multirow{2}{*}{\textbf{\texttt{change\_request\_ratio}}} 
 & \multirow{2}{*}{\pbox{4.5cm}{Dynamically scales a user's workload generation frequency.}} 
 & \texttt{multiplier} & Probability distribution generating the multiplier applied to the user's request ratio. \\
\bottomrule
\end{tabular*}

\end{table*}
\end{center}



In Edge and 5G/6G environments, user mobility and connectivity volatility are the primary drivers for workload migrations. \textsc{ECCO-Generator} models user behavior through continuous movement, request rate fluctuations, and temporary or permanent disconnections (Table~\ref{tab:user_actions})~\cite{Haibeh2022A}.

The most basic user-level behavior is permanent disconnection, represented by removing a user from the system (\texttt{remove\_user}). This event only requires modeling its statistical occurrence and does not require any additional impact model. On contrary, temporary disconnections (\texttt{suspend\_user}) requires of a more fine-grained configuration.

A user can lose connectivity either by moving outside the coverage radius, resulting in a Radio Link Failure~\cite{Ampratwum2025Radio}, or stochastically through the suspend user event, which models temporary disconnections such as tunnel traversal, device hardware failures, sudden network outages, or IoT deep-sleep cycles~\cite{Haibeh2022A}. In the latter case, the configuration must define both the statistical distribution governing how often a given user is temporarily disconnected and the impact distribution that determines the disconnection duration, that is, the time elapsed until the user is reconnected.


In the case of user mobility, the simulation framework decouples the intrinsic mobility properties of user entities from the stochastic triggers that govern their physical displacement, represented by the move user event (\texttt{remove\_user}). The temporal execution of spatial displacement is strictly driven by the move user event, whose occurrence frequency must be explicitly defined via stochastic distributions in the simulation configuration for \textsc{ECCO-Generator}, enhancing the realism and variability of the spatial dynamics~\cite{Bettstetter2004Stochastic}.


The mobility model, user speed, trajectory-change probabilities, and coverage radius are defined by \textsc{ECCO-Architect} (Section~\ref{sec:ecco_architect}). Under the Manhattan mobility model, each move user event extends the user's previous trajectory by advancing along the orthogonal street grid. When the user reaches a street intersection, the direction-change probabilities are evaluated, and any remaining movement distance is applied along the selected outgoing direction~\cite{Kim2009Realistic}.

Conversely, the Random Direction mobility model is not constrained by predefined paths. Each user preserves its previous direction vector, and every \texttt{move\_user} event first evaluates the configured probability of changing direction~\cite{Khan2016Modeling}. If no change is triggered, the user continues along the same vector; otherwise, a new direction is sampled and the position is updated according to the user's speed.


Once the user's new position is computed by the mobility model, \textsc{ECCOS} evaluates its proximity to the edge infrastructure. The user is dynamically associated with the nearest edge node when it remains within the configured coverage radius; otherwise, a signal-loss condition is triggered and the user's connectivity state is temporarily suspended~\cite{Qiu2025Mobility}. Consequently, not every move user event produces a roaming or disconnection action. A roaming action only occurs when the nearest edge node changes due to the user's new geographical position, whereas a disconnection only occurs when the user moves outside the coverage range of all available edge connection nodes~\cite{Loutfi2024An}. Users disconnected as a result of mobility remain disconnected until subsequent movement events bring them back within the coverage radius of an edge node.

Finally, the browsing behaviour of the users can be also modified during the simulation by altering the request ratio for the application that the user requests (\texttt{change\_request\_ratio}), emulating the highly dynamic nature of real-world mobile traffic~\cite{Li2023Estimating}. The change request ratio is governed by its inter-arrival time distribution. The event incorporates a directional probability threshold that model the anomaly as a traffic burst (multiplying the current request ratio by a sampled scalar) or a traffic valley (dividing it by the same scalar). The event behaves as a bounded, self-restoring state machine. Each triggered fluctuation dynamically schedules a secondary restoration event, parameterized by its own time distribution. Furthermore, the architecture guarantees state consistency under edge cases; if a user suffers a physical disconnection or hardware failure during an active burst, the system sanitizes the variables to ensure that the user resumes normal traffic behavior upon reconnection, preserving the integrity of the simulation.



\subsubsection{Application Events}
\label{sect_applicationevents}


\begin{center}
\begin{table*}[!ht]%
\caption{Stochastic action catalog and impact configuration schema for the \texttt{application} domain.\label{tab:app_actions}}
\begin{tabular*}{\textwidth}{@{\extracolsep\fill}ll l p{5.5cm}@{\extracolsep\fill}}
\toprule
\multicolumn{2}{@{}l}{\textbf{Action Specification}} & \multicolumn{2}{@{}l}{\textbf{Impact Configuration Schema}} \\\cmidrule{1-2}\cmidrule{3-4}
\textbf{Action Name} & \textbf{Action Description}  & \textbf{Impact Attribute}  & \textbf{Attribute Description \& Semantic Meaning}   \\
\midrule
\textbf{\texttt{remove\_app}} 
 & Decommissions an application graph. 
 & -- & -- \\
\midrule
\multirow{3}{*}{\textbf{\texttt{update\_app\_footprint}}} 
 & \multirow{3}{*}{\pbox{4cm}{Scales computational demands (CPU, RAM\ldots) of microservices.}} 
 & \texttt{multiplier} & Distribution scaling computational resource requirements of service nodes. \\
 & & \texttt{sigma\_footprint} & Noise distribution/standard deviation introduced during footprint resizing. \\
\midrule
\multirow{3}{*}{\textbf{\texttt{update\_app\_network}}} 
 & \multirow{3}{*}{\pbox{4cm}{Alters network links demands generated across service intercomunication edges.}} 
 & \texttt{multiplier} & Distribution scaling network requeriment demands. \\
 & & \texttt{sigma\_net} & Noise distribution introduced into service link network constraints. \\
\midrule
\multirow{3}{*}{\textbf{\texttt{update\_app\_topology}}} 
 & \multirow{3}{*}{\pbox{4cm}{Mutates service graph structure by adding or removing microservices.}} 
 & \texttt{l\_min} & Distribution specifying the lower bound for microservice count in the DAG/SFC. \\
 & & \texttt{l\_max\_global} & Distribution specifying the upper bound for microservice count in the DAG/SFC. \\
\midrule
\multirow{3}{*}{\textbf{\texttt{surge\_popularity}}} 
 & \multirow{3}{*}{\pbox{4cm}{Triggers a transient upward surge in application request volume.}} 
 & \texttt{transient\_prob} & Probability distribution determining whether the spike is transient or permanent. \\
 & & \texttt{duration\_dist} & Distribution governing the duration of a transient surge before returning to normal. \\
\midrule
\multirow{3}{*}{\textbf{\texttt{drop\_popularity}}} 
 & \multirow{3}{*}{\pbox{4cm}{Induces a transient decrease in application popularity and request load.}} 
 & \texttt{transient\_prob} & Probability distribution determining whether the drop is transient or permanent. \\
 & & \texttt{duration\_dist} & Distribution governing the duration of a transient popularity drop. \\
\midrule
\multirow{3}{*}{\textbf{\texttt{geo\_demand\_shift}}} 
 & \multirow{3}{*}{\pbox{4cm}{Relocates geographic center of gravity of application demand.}} 
 & \texttt{jump\_prob} & Probability distribution determining abrupt geographic jump vs. smooth drift. \\
 & & \texttt{drift\_dist} & Spatial displacement distribution $(\Delta X, \Delta Y)$ applied during smooth drifts. \\
\bottomrule
\end{tabular*}

\end{table*}
\end{center}



Modern applications deployed in the Continuum are not monolithic nor static; they rely on microservices composition and experience geographical and temporal demand shifts~\cite{Kadouma2025A}. Beyond simple application deletion (\texttt{remove\_app}), applications may undergo lifecycle updates that affect their resource demands or service composition (Table~\ref{tab:app_actions})~\cite{Baresi2019A}. Day-two operations can significantly alter an application's footprint. The simulator provides three mechanisms for modeling application lifecycles: changes in computational workload (\texttt{update\_app\_footprint}), changes in network demand (\texttt{update\_app\_network}), and changes in service composition(\texttt{update\_app\_topology}).

The events that simulate changes in application workloads and network usage operate in a similar manner. They respectively modify the resource requirements of constituent microservices and the application's network requirements. Their stochastic impact is parameterized by a user-defined probability distribution that represents the standard deviation of the applied noise. This perturbation either proportionally scales the application's computational resource demands during its lifecycle or simulates sudden network demand spikes and varying Quality of Service (QoS) conditions~\cite{Fu2022Adaptive}.


The event that changes an application's service composition is applicable only to microservice-based architectures. It simulates the dynamic evolution of applications through the stochastic insertion or deletion of microservices, while enforcing lower and upper bounds on the number of services per application~\cite{Fu2022Adaptive}.


\textsc{ECCOS} simulates viral trends or "Flash Crowds" through events that abruptly modify an application's popularity, either by increasing it through a surge (\texttt{surge\_popularity}) or decreasing it through a drop (\texttt{drop\_popularity})~\cite{10.1145/2068816.2068867}. This modification only affects the application-selection process of users created after the event, thereby increasing or decreasing the number of new users likely to request the application. Existing application popularity values and the request ratios of users already associated with these applications remain unchanged, and they are altered only by user's events. The phenomenon can be transient or permanent. When it is transient, as in the case of a live sporting event, a paired restore-popularity event is scheduled after a specified period of time~\cite{10.1145/2068816.2068867}.

Finally, and for the cases of localized applications, defined as subsets of applications with higher popularity in hot-spot areas, the emulator allows to simulate the spatial migration of user interest (\texttt{geo\_demand\_shift}). This is achieved by displacing the hot spot either through random jumps to new coordinates or through small incremental movements from its current location. As a result, the optimization engine may need to perform live migrations of edge microservices to track the evolving user base~\cite{Wang2019Delay}.




\subsubsection{Global Spawner Events}
\label{subsec:global_spawner}


\begin{center}
\begin{table*}[!ht]%
\caption{Stochastic action catalog and impact configuration schema for the \texttt{global\_spawner} domain.\label{tab:global_actions}}
\begin{tabular*}{\textwidth}{@{\extracolsep\fill}ll l p{6.8cm}@{\extracolsep\fill}}
\toprule
\multicolumn{2}{@{}l}{\textbf{Action Specification}} & \multicolumn{2}{@{}l}{\textbf{Impact Configuration Schema}} \\\cmidrule{1-2}\cmidrule{3-4}
\textbf{Action Name} & \textbf{Action Description}  & \textbf{Impact Attribute}  & \textbf{Attribute Description \& Semantic Meaning}   \\
\midrule
\multirow{1}{*}{\textbf{\texttt{new\_user}}} 
 & \multirow{1}{*}{\pbox{4.5cm}{Dynamically instantiates a new user from scratch (\textit{ex novo}).}} 
 & -- & -- \\
  & 
 &  &  \\
\midrule
\multirow{2}{*}{\textbf{\texttt{new\_app}}} 
 & \multirow{2}{*}{\pbox{4.5cm}{Dynamically instantiates a new application DAG/SFC from scratch.}} 
 & \texttt{num\_new\_users} & Distribution determining the number of initial users spawned for the new app. \\
\bottomrule
\end{tabular*}
\end{table*}
\end{center}



Creating new users (\texttt{new\_user}) and applications (\texttt{new\_app}) is a global simulation event that represents the external arrival of a new entity into the scenario. The global spawner is the \textsc{ECCO-Generator} component responsible for introducing these new simulation entities (Table~\ref{tab:global_actions})~\cite{Schriber2004Inside}. Its role differs from actions attached to existing instances: user or application actions modify, move, remove, or update specific objects, whereas creation events add entirely new objects to the simulation state. This distinction is necessary because the internal event-generation logic assigns configured actions individually to each existing entity. For example, if a remove-user event is configured, the simulator schedules an independent removal countdown for User 1, another for User 2, and so on. Such behavior is appropriate for instance-specific events, but not for entity creation, which must instead be defined in a general event container such as the global spawner~\cite{Ullrich2017An}.

This design choice preserves the intended DES semantics: one global arrival process controls the birth of new entities, while per-object action blocks control the lifecycle and behavior of entities after they have been instantiated. This design cleanly separates exogenous simulation dynamics from endogenous entity behavior~\cite{Ullrich2017An}.


\subsubsection{User Composition of Encapsulated Events}
\label{sect_usercompositionofencapsulatedevents}

A significant challenge in discrete-event simulation for edge and cloud environments is capturing the complex, cascading nature of real-world phenomena~\cite{Higashino2016CEPSim}. While traditional simulators treat events as isolated occurrences affecting a single entity (e.g., a node failure, a user arrival), real-world occurrences often trigger a chain of heterogeneous impacts across the infrastructure. For instance, a catastrophic power outage not only disables a physical node but simultaneously terminates the applications hosted within it and disconnects the locally attached users. Furthermore, macroscopic events such as localized flash crowds require coordinated actions across multiple entities. 

\textsc{ECCOS} simulators allows user-defined encapsulated events, which acts as a logical container that defines a single probability distribution for its occurrence frequency, but its execution impact is decomposed into a collection of primitive sub-actions affecting heterogeneous entities~\cite{Zhu2017A}. 

This is achieved with a declarative encapsulated events architecture, and defining the use of the \texttt{composed\_of} directive within the event's YAML configuration. When the simulator's event engine dispatches the macro-event, it iterates through the \texttt{composed\_of} list, dynamically resolving dependencies and executing the underlying base actions within the same simulation time step. 

The simulator's compositional event architecture is designed to be completely orthogonal to the origin of the trigger. Consequently, composite events can be associated both with concrete domain instances or with the global spawner to orchestrate multi-step structural dynamics.

When an encapsulated event involves multiple entities, identifying the specific instance over which the sub-action must be applied is a critical challenge. To solve this, the simulator implements a dynamic target resolution engine. This engine provides granular control over how target instances are dynamically resolved during the simulation loop. The resolution engine offers several configuration alternatives to identify the affected instances: the same instance that generates the macro-event; a uniform random selection from all the instances; an specific instance idintified by its identificator; a delegated and self-resolution of the simulator; or to a group of instances.

In the latter case, when the dynamic target-resolution engine is configured to select a group of instances, the event affects a cluster of $N$ elements simultaneously~\cite{Roumeliotis2025Multi}. The cluster size can be defined either as an absolute integer (e.g., exactly 5 nodes) or dynamically as a probability or percentage (e.g., $0.15$ to affect 15\% of the currently active users in the simulator). The target resolution of this group requires additional parametrization, supporting explicit list selection~\cite{Xiong2019A}, uniform random sampling~\cite{Roumeliotis2025Multi}, a hybrid approach that selects the originating instance plus $N-1$ random instances~\cite{Wu2020Bilateral}, or a proximity-based selection model~\cite{Li2020Privacy}. This last alternative considers two cases: selecting a random anchor element and its proximate elements, or selecting the \textit{self}-generating element and its proximate elements.

This contextual proximity model is particularly useful for modeling localized phenomena. Because proximity depends on the entity type, the resolution engine implements heterogeneous proximity metrics: for infrastructure nodes, proximity is defined as the topological shortest-path distance within the network graph~\cite{Li2020Privacy}; for users, it corresponds to the topological distance between the access points to which they are actively connected~\cite{Wang2020Collaborative}; and for applications, it is based either on geographical affinity, when applications share the same spatial hot-spot mapping, or on popularity adjacency, grouping applications with similar popularity ranks~\cite{Huang2026Edge}.



\subsection{ECCO-State: System State Manager}
\label{subsec:ecco_state}

The \textsc{ECCO-State} module constitutes the state persistence and analytical accounting core of the simulator, acting as the authoritative \emph{Single Source of Truth} regarding the instantaneous physical and logical status of the underlying infrastructure at any discrete time step $t$. Its primary architectural responsibility is to enforce mathematical consistency and structural invariants across the discrete-event stochastic transitions driven by \textsc{ECCO-Generator} and the placement decisions computed by \textsc{ECCO-Solver}.


At the operational level, \textsc{ECCO-State} maintains the unified structural graph of the tiered physical infrastructure $\mathcal{G}$, comprising heterogeneous compute nodes $\mathcal{N}$ and interconnecting network links $\mathcal{L}$, alongside the Directed Acyclic Graphs (DAGs) representing applications $\mathcal{A}$ and the active user population $\mathcal{U}$. Each time an event from the event list generated by \textsc{ECCO-Generator} is processed, the simulator state changes and all these structures and state variables must be updated consistently~\cite{Appuhamilage2025A}, handled by \textsc{ECCO-State}.


To guarantee physical feasibility throughout the simulation horizon, the module also implements a rigorous residual-capacity accounting engine. A fundamental architectural strength of \textsc{ECCO-State} is its dynamic, schema-agnostic resource accounting model. Unlike domain-specific edge simulators that hardcode a rigid tuple of physical capacities (typically limited to CPU cores, RAM memory, and link bandwidth), \textsc{ECCO-State} operates over arbitrary, user-defined multi-dimensional attribute spaces defined via configuration manifests~\cite{Reidys2025Coach}.


By employing dynamic introspection over graph node and edge dictionaries, the state manager natively tracks any set of custom quantitative attributes (features) $F = \{f_1, f_2, \dots, f_k\}$ specified by the user (e.g., hardware accelerators such as GPUs/TPUs, disk storage constraints, or secure enclave quotas)~\cite{Qiu2025Cloud}. During state validation and solver formulation, the system iterates over all numeric capacities present in the infrastructure topology and automatically enforces residual capacity invariants across matching microservice demands:
\begin{equation}
    \sum_{m \in \mathcal{M}_n} d_{a,m}^f \le C_n^f, \quad \forall n \in \mathcal{N}, \; \forall f \in F
\end{equation}
where $d_{a,m}^f$ represents the demand of microservice $m$ for attribute $f$, $C_n^f$ denotes the physical capacity of node $n$ for attribute $f$, and $\mathcal{M}_n$ represents the set of microservices from any application that are currently placed in node $n \in \mathcal{N}$ and determined by the current placement dictionary $\mathcal{P}^{(t)}$. If an attribute is omitted from either the infrastructure or application specification, the accounting engine safely evaluates demand to zero without schema violations, ensuring complete backward compatibility and high extensibility~\cite{Zakarya2019Modelling}.


In summary, every structural state transition, whether triggered by external stochastic events (e.g., user arrivals, mobility, or infrastructure degradation) or internal solver allocations, is transactionally validated by \textsc{ECCO-State} to prevent resource oversubscription~\cite{Qiu2025Cloud}. Furthermore, the module constructs immutable, deterministic state snapshots ($S^{(t)}$) that are exposed via standardized interfaces to optimization engines.

Finally, a distinguishing architectural capability of \textsc{ECCO-State} is its support for \emph{latent or non-linear state transitions} within application and user lifecycles, overcoming the limitations of traditional binary simulators that restrict services to mutually exclusive \emph{active} or \emph{terminated} states.

Rather than naïvely destroying deployed microservices and releasing allocated memory upon a temporary user disconnection, \textsc{ECCO-State} transitions the affected session into a \emph{latent soft-state retention mode}. In those cases, the effective request generation rate of the disconnected user drops to zero, immediately relieving network consumptions but retaining the node resources allocation, emulating static reservation and dynamic load over the resource attributes~\cite{Yi2024Mobile}.

\subsection{Ecco-Solver: Decoupled Optimization Engine}
\label{sec:ecco_solver}

A cornerstone of the \textsc{ECCOS} simulation framework is its strict architectural separation between discrete-event simulation dynamics and mathematical optimization~\cite{Soykan2024Optimizing}. The \textsc{ECCO-Solver} module acts as a completely decoupled optimization engine responsible for computing the spatial mapping of application microservices across the underlying physical infrastructure. By isolating the placement logic behind a uniform abstraction layer, \textsc{ECCOS} allows researchers to design, prototype, and benchmark custom placement algorithms, ranging from exact Integer Linear Programming (ILP) formulations to fast heuristics or Reinforcement Learning policies, without modifying the core event scheduler or system state tracking.

Every solver in ECCOS inherits from the abstract base class \texttt{BaseSolver}, which establishes a stateless, standardized interface for decision-making. When a placement policy is triggered at simulation time $t$, the simulation runner invokes the solver through a single contract:

\begin{equation}
(\mathcal{P}^{(t)}, \, \mathcal{C}^{(t)}) \leftarrow \texttt{solve}\Big(\mathcal{G}, \, \mathcal{A}, \, \mathcal{U}, \, \Omega_{\text{cfg}}, \, \mathcal{P}^{(t-\Delta t)}\Big)
\end{equation}

To implement a new custom solver (such as an exact mathematical model or a fast greedy heuristic), a developer must consider the five structured inputs provided by the simulator and the required output specification~\cite{Santos2021Service}:

\begin{enumerate}
    \item \textbf{Infrastructure Graph ($\mathcal{G}$):} A directed graph representing the active physical infrastructure. Nodes $n \in \mathcal{N}$ expose multi-dimensional resource capacities $C_n^f$ for each hardware attribute $f \in F$, such as $f \in \{\text{CPU, RAM, Storage}\}$, along with binary operational states. Edges representing network links provide all-pairs shortest-path propagation delays $D(i, j)$ for any two physical nodes $i, j \in \mathcal{N}$ and, similarly to the nodes, their multi-dimensional resource attributes, such as $f \in \{\text{Delay, Bandwidth}\}$.
    \item \textbf{Application Workload ($\mathcal{A}$):} The active set of Service Function Chaining (SFC) or Directed Acyclic Graph (DAG) applications. Each application $a \in \mathcal{A}$ consists of a set of ordered microservices $m \in \mathcal{M}_a$ with computational demands $d_{a,m}^f$, connected by directed virtual links that models the microservice's interactions in application $a$ $\mathcal{I}_a = \{(m_1, m_2)\}$.
    \item \textbf{User Population ($\mathcal{U}$):} The active mobile user instances generating traffic requests. Each user $u \in \mathcal{U}$ is characterized by their requested application $a(u) \in \mathcal{A}$, their currently associated network access point $n(u) \in \mathcal{N}$, and their dynamic request ratio multiplier $\lambda_u$.
    \item \textbf{Execution Configuration ($\Omega_{\text{cfg}}$):} Runtime hyperparameters loaded from the solver YAML configuration file. This parameters are solver-specific, and different parameters would be consider for different solvers. For example, for ILP solvers, the list of parameters could include optimality tolerances (\texttt{gapRel}), execution timeouts (\texttt{timeLimit}), and objective weighting coefficients.
    \item \textbf{Previous System State ($\mathcal{P}^{(t-\Delta t)}$):} A structured snapshot of the placement dictionary at the preceding decision epoch. Providing the historical placement is critical for stateful optimization, enabling solvers to evaluate service migration costs and prevent hysteresis or route flapping.
\end{enumerate}


Regardless of internal algorithmic complexity, every solver must return a standardized tuple $(\mathcal{P}, \mathcal{C})$. To accommodate both single-instance and multi-replica distributed deployments across the Computing Continuum~\cite{Santos2021Service}, the placement mapping $\mathcal{P}$ is formulated as a nested dictionary assigning each microservice unique identifier $m \in \mathcal{M}_a$ of application $a \in \mathcal{A}$ to a non-empty subset of physical hosting nodes $\mathcal{N}_{a,m} \subseteq \mathcal{N}$:
\begin{equation}
\mathcal{P} = \Big\{ a: \big\{ m: \mathcal{N}_{a,m}, \; \forall m \in \mathcal{M}_a \big\}, \; \forall a \in \mathcal{A} \Big\}
\end{equation}
In standard non-replicated deployments, the hosting set contains a single target node ($|\mathcal{N}_{a,m}| = 1$), whereas in multi-instance load-balancing or redundant workloads, $\mathcal{N}_{a,m}$ contains multiple simultaneous host nodes ($|\mathcal{N}_{a,m}| \geq 1$). The scalar $\mathcal{C} \in \mathbb{R}^+$ represents the evaluated objective cost of the resulting placement mapping, or returns a heavy penalty ($\mathcal{C} = \texttt{infeasible\_penalty}$) if the problem is mathematically infeasible due to physical capacity, or other, constraints.


In a high-frequency Discrete-Event Simulation where hundreds of stochastic events occur per simulated minute, invoking an NP-hard mathematical solver on every single event would cause computational bottlenecks~\cite{He2024An}. To ensure simulation scalability, \textsc{ECCOS} incorporates a dedicated trigger policy manager that decouples event arrival frequency from solver execution frequency or conditions~\cite{Peng2022Distributed}.

Configurable via the solver YAML specification attribute \texttt{trigger\_policy.type} (Section~\ref{append:solvertriggerpolicy}), the trigger manager evaluates every incoming simulation event against a user-defined policy to decide whether to trigger solver recalculation:
\begin{itemize}
    \item \textbf{Every Event (\texttt{solve\_all}):} Triggers the solver after each single event.
    \item \textbf{Solver Disable (\texttt{solve\_none}):} The solver is not executed after any event. The output only includes the list of events in the system without the placement solutions.
    \item \textbf{Periodic Time Intervals (\texttt{solve\_every\_t\_seconds}):} Invokes the solver strictly at fixed simulation time periods \texttt{interval\_seconds}.
    \item \textbf{Event Batching (\texttt{solve\_every\_n\_events}):} Triggers recalculation after every \texttt{batch\_size} cumulative discrete events across the ecosystem.
    \item \textbf{User-listed Event Interrupts (\texttt{solve\_on\_event\_types}):} An event-driven policy that executes the solver only upon the arrival of a limited list of events in attribute \texttt{critical\_events}.
    \item \textbf{Temporal Windows (\texttt{solve\_time\_windows}):} Confines optimization execution to specific lists (\texttt{windows}) of simulation time slots $[t_{\text{start}}, t_{\text{end}}]$.
    \item \textbf{Event Index Ranges (\texttt{solve\_time\_windows}):} Confines optimization execution to specific lists (\texttt{windows}) of events intervals $[i_{\text{start}}, i_{\text{end}}]$.
    \item \textbf{Random Event Selection (\texttt{solve\_random\_prob}):} Triggers solver execution probabilistically for each event according to the specified \texttt{probability} value.
    \item \textbf{Pattern-based Selection (\texttt{solve\_custom\_pattern}):} A sequence of Boolean values manually specifies the events for which the solver must be triggered (\texttt{pattern}). This pattern is applied cyclically throughout the simulation.
    \item \textbf{Hybrid Policy (\texttt{combined}):} Combines periodic time-driven recalculations every \texttt{interval\_seconds} seconds (\texttt{solve\_every\_t\_seconds}), with immediate event-driven interrupts whenever a user-listed event in \texttt{critical\_events} (\texttt{solve\_on\_event\_types}).
\end{itemize}


Note that \texttt{type:solve\_none} uses only the event-trace generation capabilities of \textsc{ECCOS}, but it opens the possibility of employing the framework as a discrete-event generator that can be integrated with third-party applications lacking an equally powerful mechanism for defining and producing discrete events (Section XXXXXX).

\textsc{ECCOS} currently provides three built-in solvers: two ILP-based formulations (Sections~\ref{subsec:ilp_single} and~\ref{subsec:ilp_multi}) and one greedy heuristic (Section~\ref{subsec:greedy_solver}). The solver used in a given simulation is selected in the solver YAML configuration file through the \texttt{setup.solver} attribute. In any case, as discussed above, users can extend the set of available solvers by implementing the \texttt{BaseSolver} interface, and then make the new solver visible and selectable from the YAML configuration by registering it through the \texttt{SolverFactory} class.


% =========================================================================
\subsubsection{Single-Objective ILP: Pure Latency Optimization}
\label{subsec:ilp_single}

The first baseline formulation implemented in ECCOS (\texttt{ILPSingleObjectiveSolver}) models the classic latency-minimal service  placement problem. Its sole objective is to minimize the total weighted end-to-end network delay perceived by active users, accounting for both ingress access latency and internal microservice-to-microservice communication delay~\cite{Soumplis2024Performance}.

\textbf{Decision Variables.} The formulation defines two sets of decision variables:
\begin{align}
x_{a,m,n} &\in \{0, 1\}, && \forall a \in \mathcal{A}, \; m \in \mathcal{M}_a, \; n \in \mathcal{N} \\
y_{a,i,n_1,n_2} &\in [0, 1], && \forall a \in \mathcal{A}, \; i=(m_1,m_2) \in \mathcal{I}_a, \; n_1, n_2 \in \mathcal{N}
\end{align}
where $x_{a,m,n}=1$ if microservice $m$ of application $a$ is hosted on physical server $n$. The continuous variable $y_{a,e,n_1,n_2}$ represents the linearized link placement, taking value $1$ if virtual interaction edge $i=(m_1, m_2)$ is mapped across the physical path from $n_1$ to $n_2$~\cite{Taghavian2023An}.

\textbf{Objective Function.} Let $m_0^{(a)}$ denote the ingress (first) microservice of application $a$, and let $\Lambda_a = \sum_{u \in \mathcal{U}: a(u)=a} \lambda_u$ represent the aggregated request load for application $a$~\cite{Soumplis2024Performance}. The objective function minimizes total weighted latency:
\begin{equation}
\min \quad \sum_{u \in \mathcal{U}} \lambda_u \sum_{n \in \mathcal{N}} D(n(u), \, n) \cdot x_{a(u), \, m_0^{(a(u))}, \, n}
+ \sum_{a \in \mathcal{A}} \Lambda_a \sum_{i=(m_1,m_2) \in \mathcal{I}_a} \sum_{n_1 \in \mathcal{N}} \sum_{n_2 \in \mathcal{N}} D(n_1, n_2) \cdot y_{a,i,n_1,n_2}
\end{equation}

\textbf{Constraints.} The optimization is subject to exact placement, physical capacity, and flow linearization constraints~\cite{Taghavian2023An}:
\begin{align}
\sum_{n \in \mathcal{N}} x_{a,m,n} &= 1, && \forall a \in \mathcal{A}, \; m \in \mathcal{M}_a \label{eq:single_exact} \\
\sum_{a \in \mathcal{A}} \sum_{m \in \mathcal{M}_a} d_{a,m}^k \cdot x_{a,m,n} &\leq C_n^k, && \forall n \in \mathcal{N}, \; k \in \{\text{CPU, RAM}\} \label{eq:single_cap} \\
\sum_{n_2 \in \mathcal{N}} y_{a,i,n_1,n_2} &= x_{a,m_1,n_1}, && \forall a \in \mathcal{A}, \; i=(m_1,m_2), \; n_1 \in \mathcal{N} \label{eq:lin1} \\
\sum_{n_1 \in \mathcal{N}} y_{a,i,n_1,n_2} &= x_{a,m_2,n_2}, && \forall a \in \mathcal{A}, \; i=(m_1,m_2), \; n_2 \in \mathcal{N} \label{eq:lin2}
\end{align}
Constraint~\eqref{eq:single_exact} guarantees that every microservice instance is deployed on exactly one operational server. Constraint~\eqref{eq:single_cap} ensures that aggregate resource consumption does not exceed node capacity. Constraints~\eqref{eq:lin1} and~\eqref{eq:lin2} enforce topological flow conservation between microservice placement variables and virtual link variables. These two flow linearization constraints~\eqref{eq:lin1} and~\eqref{eq:lin2} are required to avoid quadratic non-linearities arising from the multiplication of binary placement variables, and the formulation introduces auxiliary link-mapping variables $y_{a,i,n_1,n_2} \in {0,1}$. Specifically, these constraints  enforce that an application edge $i=(m_1,m_2)$ is mapped over physical path $(n_1,n_2)$ if and only if microservices $m_1$ and $m_2$ are hosted on servers $n_1$ and $n_2$, respectively. This guarantees exact accounting of inter-service network delays while preserving a computationally tractable Mixed-Integer Linear Programming (MILP) formulation.

% =========================================================================
\subsubsection{Multi-Objective ILP: Latency, Migration, and Consolidation}
\label{subsec:ilp_multi}

While pure latency optimization achieves optimal quality of service, applying it independently at each time step often leads to operational instability: minor shifts in user mobility can trigger massive microservice relocations across servers. To address stateful trade-offs, \textsc{ECCOS} implements a second solver (\texttt{ILPMultiObjectiveSolver}) that balances end-to-end latency against service migration overhead and server consolidation (energy efficiency).

\textbf{Decision Variables.} In addition to $x_{a,m,n}$ and $y_{a,e,n_1,n_2}$, the multi-objective model introduces a binary server activation variable:
\begin{equation}
z_n \in \{0, 1\}, \quad \forall n \in \mathcal{N}
\end{equation}
where $z_n = 1$ if server node $n$ hosts at least one active microservice instance.

\textbf{Multi-Objective Cost Formulation.} Let $\mathcal{P}^{(t-\Delta t)}(a,m) \in \mathcal{N}$ denote the physical server that hosted microservice $m$ of application $a$ in the previous simulation epoch. The migration penalty indicator is formulated as $(1 - x_{a,m,\mathcal{P}^{(t-\Delta t)}(a,m)})$, which evaluates to $1$ whenever a microservice is moved to a different node.

The composite scalar objective function minimizes a weighted linear combination of latency, migration frequency, and active server footprint:
\begin{equation}
\min \quad w_{\text{lat}} \cdot \Phi_{\text{lat}} + w_{\text{mig}} \cdot \Phi_{\text{mig}} + w_{\text{srv}} \cdot \Phi_{\text{srv}}
\end{equation}
where the individual cost components are defined as:
\begin{align}
\Phi_{\text{lat}} &= \sum_{u \in \mathcal{U}} \lambda_u \sum_{n \in \mathcal{N}} D(u(n), n) x_{a,m_0,n} + \sum_{a,i,n_1,n_2} \Lambda_a D(n_1,n_2) y_{a,i,n_1,n_2} \\
\Phi_{\text{mig}} &= \sum_{a \in \mathcal{A}} \sum_{m \in \mathcal{M}_a} \mathbb{D}_{\text{active}}(a,m) \cdot \Big( 1 - x_{a, \, m, \, \mathcal{P}^{(t-\Delta t)}(a,m)} \Big) \\
\Phi_{\text{srv}} &= \sum_{n \in \mathcal{N}} z_n
\end{align}
where the active deployment indicator function $\mathbb{D}_{\text{active}}(a,m) \in {0, 1}$ is evaluated to 1 if microservice $m \in \mathcal{M}_a$ was actively running in the preceding time step $(t-\Delta t)$, ensuring that newly instantiated microservices do not trigger false migration penalties.
The non-negative weights $w_{\text{lat}}$, $w_{\text{mig}}$, and $w_{\text{srv}}$ are configurable via the solver configuration YAML file, allowing users to tune the simulator between aggressive low-latency placement and highly stable, consolidated deployments.

\textbf{Additional Constraints.} In addition to constraints~\eqref{eq:single_exact}--\eqref{eq:lin2}, the multi-objective solver couples microservice assignments to server activation status via a Big-$M$ constraint~\cite{9060940}:
\begin{equation}
\sum_{a \in \mathcal{A}} \sum_{m \in \mathcal{M}_a} x_{a,m,n} \leq M \cdot z_n, \quad \forall n \in \mathcal{N}
\end{equation}
where $M = \sum_{a \in \mathcal{A}} |\mathcal{M}_a|$ represents an upper bound on the total number of deployable microservices across the entire ecosystem.

% =========================================================================
% =========================================================================
\subsubsection{Demand-Aware Greedy Heuristic: Sub-Millisecond Attraction Placement}
\label{subsec:greedy_solver}

Exact Integer Linear Programming solvers guarantee global or Pareto-optimal placements but suffer from exponential worst-case computational complexity ($O(2^{|\mathcal{A}||\mathcal{M}||\mathcal{N}|})$). As infrastructure topologies scale to hundreds of Edge and Fog nodes, solving ILPs at every decision epoch can introduce undesirable simulation bottlenecks. To provide an ultra-fast baseline suitable for large-scale real-time orchestration, \textsc{ECCOS} includes a deterministic heuristic solver (\texttt{GreedySolver}) based on a demand-aware topological attraction strategy.

\textbf{Application Prioritization.} Rather than jointly optimizing all applications simultaneously, the greedy solver decomposes the placement problem sequentially by prioritizing applications with the highest user demand. Let $\mathcal{U}_a = \{u \in \mathcal{U} : a(u) = a\}$ denote the subset of active users requesting application $a \in \mathcal{A}$. The aggregate request intensity $\Lambda_a$ for application $a$ is formulated as:
\begin{equation}
\Lambda_a = \sum_{u \in \mathcal{U}_a} \lambda_u
\end{equation}
The solver sorts the application set $\mathcal{A}$ into an ordered sequence $\vec{\mathcal{A}} = (a_{(1)}, a_{(2)}, \dots, a_{(|\mathcal{A}|)})$ such that $\Lambda_{a_{(1)}} \geq \Lambda_{a_{(2)}} \geq \dots \geq \Lambda_{a_{(|\mathcal{A}|)}}$, breaking ties deterministically by unique application identifier.

\textbf{Topological Attraction Scoring.} For each prioritized application $a \in \vec{\mathcal{A}}$, the heuristic evaluates all operational server nodes $n \in \mathcal{N}$ by computing a topological attraction score $\Psi(a, n)$. This metric quantifies the cumulative weighted shortest-path access latency from all users currently requesting application $a$ to candidate server $n$:
\begin{equation}
\Psi(a, n) = \sum_{u \in \mathcal{U}_a} \lambda_u \cdot D(u_{\text{node}}, \, n), \quad \forall n \in \mathcal{N}
\end{equation}
Candidate physical servers are ranked in ascending order of $\Psi(a, n)$, producing a prioritized node list where servers topologically closest to the application's active user centroid appear first.

\textbf{Monolithic-First Allocation Strategy.} Let $R_n^k(t)$ denote the remaining available capacity of resource $k \in \{\text{CPU, RAM}\}$ on physical node $n$ at the current step of the greedy allocation loop. For each application $a$, the placement mapping is constructed via a two-stage allocation mechanism:
\begin{enumerate}
    \item \textbf{Monolithic Co-location Attempt:} The solver first scans the ranked candidate nodes to find the best server $n^*$ capable of co-locating the entire microservice chain $\mathcal{M}_a$ locally, satisfying:
    \begin{equation}
    \sum_{m \in \mathcal{M}_a} d_{a,m}^k \leq R_{n^*}^k(t), \quad \forall k \in \{\text{CPU, RAM}\}
    \end{equation}
    If such a node exists, all microservices $m \in \mathcal{M}_a$ are assigned to $n^*$ ($\mathcal{P}_{a,m} \leftarrow n^*$), eliminating inter-microservice network delay ($D(n^*, n^*) = 0$). Remaining node capacities are immediately decremented: $R_{n^*}^k(t) \leftarrow R_{n^*}^k(t) - \sum_m d_{a,m}^k$.
    \item \textbf{Distributed Fallback Allocation:} If no single server possesses sufficient residual resources to host the entire application, the solver iterates through each microservice $m \in \mathcal{M}_a$ sequentially and deploys it on the first candidate server $n \in \mathcal{N}$ in the attraction-ranked list satisfying individual microservice demand ($d_{a,m}^k \leq R_n^k(t)$).
\end{enumerate}

\textbf{Feasibility and Output Evaluation.} If any microservice $m \in \mathcal{M}_a$ cannot be allocated due to infrastructure resource exhaustion, the algorithm terminates immediately and returns an infeasible status tuple $(\emptyset, \, \text{infeasible\_penalty})$. Otherwise, upon successfully placing all applications, the solver evaluates the standardized objective cost $\mathcal{C} \in \mathbb{R}^+$ across the finalized mapping $\mathcal{P}$ using the identical end-to-end weighted latency function applied by the single-objective ILP model, ensuring fair cross-algorithm benchmarking.







\subsection{Ecco-Bridge: Data Ecosystem and ML-Interface Engine}
\label{subsec:eccobridge}

\textsc{ECCOS} exports highly structured, self-contained trace records. Rather than dumping unstructured log streams, the simulation engine formalizes every orchestration step into an independent, well-typed JSON snapshot designed to directly support data-driven modeling adapted to Machine Learning (ML) and Reinforcement Learning (RL) post-processing.

While tabular formats such as CSV or streaming formats such as JSON Lines are standard for flat, purely numerical datasets, they prove fundamentally inadequate for representing the topological and hierarchical complexity of the Computing Continuum. A Computing Continuum state inherently comprises a non-Euclidean graph topology consisting of heterogeneous physical nodes, dynamic network edges with fluctuating inter-node propagation delays, and multi-layered service dependencies (users connected to applications composed of microservices deployed on hardware nodes). Flattening such nested graph structures into single-line CSV rows or JSON Lines strings inevitably destroys spatial adjacency and forces artificial serialization choices~\cite{Peng2025Unifying}. On contrary, hierarchical JSON provides an intuitive, human-readable representation of graph attributes (\texttt{node\_information}, \texttt{edges}, user-to-node connectivity, and application mappings). Modern ML pipelines, specifically GNN frameworks such as \emph{PyTorch Geometric} or \emph{DGL}, consume structured JSON records natively, converting each file directly into structured graph tensors without loss of relational fidelity.


Each simulation run generates a structured trace repository containing individual JSON snapshot files for each simulation step, i.e., for each event processing. An important design principle of our output architecture is that each step file is fully self-contained~\cite{Zhang2023Efficient}. It encapsulates the physical pre-decision state of the Computing Continuum ($\mathcal{G}^{(t-\Delta t)} \cup \mathcal{A}^{(t-\Delta t)} \cup \mathcal{U}^{(t-\Delta t)} \cup \mathcal{P}^{(t-\Delta t)}$), the quantitative optimization metrics achieved by the solver at the preceding decision epoch ($\mathcal{O}^{(t-\Delta t)}$), the exogenous environmental event that triggers the simulation step ($\mathcal{E}^{(t)}$), the explicit placement differential or migration plan  obtained from the difference between the previous placement and the new  placement decision obtained by the solver in the current epoch ($\Delta \mathcal{P}^{(t)} = \mathcal{P}^{(t)} \ominus \mathcal{P}^{(t-\Delta t)}$), the resulting post-decision state ($\mathcal{G}^{(t)} \cup \mathcal{A}^{(t)} \cup \mathcal{U}^{(t)} \cup \mathcal{P}^{(t)}$), and the new quantitative optimization metrics achieved by the solver ($\mathcal{O}^{(t)}$), such as end-to-end latency and resource occupancy.

Storing complete pre- and post-transition snapshots in individual files provides significant operational flexibility for dataset curation. In particular, it enables researchers to randomly sample, partition, or shuffle individual trace files across large simulation runs, whether for stochastic mini-batch training in deep imitation learning or for constructing offline replay buffers, without needing to parse or reconstruct sequential dependencies from a monolithic log stream.


Because these traces retain the full simulation context, the resulting ML and RL tasks can be formulated from multiple perspectives. This flexibility allows the same trace repository to support different learning formulations without regenerating the underlying simulation data. For example, one specific way of exploiting this flexibility is to interpret each JSON record as a Markov Decision Process (MDP) transition tuple~\cite{Sheng2021Deep}:
\begin{equation}
    \mathcal{T}_k = (S_k, A_k, R_k, S_{k+1})
\end{equation}
Under this perspective, the JSON fields can be associated with the elements of an RL or imitation-learning problem as summarized in Table~\ref{tab:rl_mapping} and detailed below~\cite{Liu2023Task}:
\begin{itemize}
    \item \textbf{Decision State ($S_k$):} The state observed by the learning agent at step $k$ can be defined as the composite tuple $S_k = (\mathcal{G}_k, \mathcal{A}_k, \mathcal{U}_k, \mathcal{P}_k, \mathcal{E}_k)$, where immediately prior to the decision state is captured in terms of the scenarion state (respectively, the physical infrastructure, the applications, the users, and the placement decison) and the exogenous system event ($\mathcal{E}_k$) that triggered the orchestration cycle.
    \item \textbf{Learning Action ($A_k$):} In this particular mapping, the action can be defined as the orchestration response to the observed state and event. For example, $A_k$ may correspond to the placement or migration plan encoded by ($\Delta \mathcal{P}^{(t)}$). In an imitation-learning setting, the same fields can instead be interpreted as the ground-truth label ($Y_k$) produced by the reference solver.
    \item \textbf{Future State ($S_{k+1}$):} The next state corresponds to the continuum configuration after applying the selected orchestration action, and is reconstructed from the post-execution dictionaries.
\end{itemize}

But as it was mentioned, other mappings are also possible. For instance, a learning task may treat the simulator-generated events ($\mathcal{E}_k$) themselves as the action sequence to be predicted (($A_k$))~\cite{Prudencio2022A}.

The reward $R_k$ of the transition tuple $\mathcal{T}_k$ could be calculated as a mathematical function of the quantitative pre- and post-optimization metrics ($\mathcal{O}^{(t-\Delta t)}$, $\mathcal{O}^{(t)}$). By exposing the granular pre- and post-decision physical metrics rather than a prescriptive scalar reward, the trace architecture remains entirely agnostic to the ML designer's reward shaping. Researchers can freely formulate, weight, or constrain custom reward functions:
\begin{equation}
    R_k = f\left(\Delta\text{o}_1, ..., \Delta\text{o}_{i}\right), o_i \in \mathcal{O}^{(t-\Delta t)}, \mathcal{O}^{(t)}
\end{equation}
during training without altering or re-simulating the underlying trace dataset. Note that his flexibility in reward definition is enabled by a distinctive architectural choice in our trace representation by the intentional omission of a hardcoded scalar reward value ($R_k$) within the exported JSON records. Instead of prescribing a single reward signal, the traces expose the raw, multi-objective telemetry captured after action execution, which is calculated by the solver~\cite{Mannion2018Reward}.

This decoupling is essential to preserve the generality and reusability of the dataset across diverse ML research paradigms~\cite{Mu2025Preference}. In real-world Computing Continuum orchestration, the definition of optimality is highly objective-dependent: one practitioner may seek to strictly minimize network latency, another may prioritize energy conservation and RAM footprint, while another may introduce penalties for service migration overhead. 

% Required packages in preamble: \usepackage{booktabs} \usepackage{tabularx}
\begin{table*}[t]
\centering
\caption{Mapping of Formal MDP Transition Elements to Simulator JSON Trace Fields}
\label{tab:rl_mapping}
\begin{tabularx}{\textwidth}{@{}l l l X@{}}
\toprule
\textbf{RL Concept} & \textbf{Output} & \textbf{JSON Labels Mapping} & \textbf{Engineering \&} \\
\textbf{(MDP Element)}   & \textbf{Trace} & & \textbf{ML Rationale} \\
\midrule
\textbf{Observation /} & $\mathcal{U}_k$ & \texttt{users\_before} dictionary & \multirow{5}{*}{\pbox{4.2cm}{Combines pre-decision topology graph attributes with the exogenous event triggering orchestration.}} \\
\textbf{State} ($S_k$) &  $\mathcal{A}_k$ & \texttt{apps\_before} dictionary& \\
&  $\mathcal{U}_k$ & \texttt{users\_before} dictionary& \\
&  $\mathcal{G}_k$ & \texttt{node\_before}/\texttt{edge\_before} dictionaries& \\
&  $\mathcal{P}_k$ & \texttt{placement\_before} dictionary& \\
&  $\mathcal{E}_k$ & \texttt{action} dictionary& \\

\addlinespace
\textbf{Agent Action} ($A_k$) & $\Delta \mathcal{P}_{k+1}$ & \texttt{diff\_message} dictionary & \multirow{3}{*}{\pbox{4.2cm}{Represents optimal service-to-node deployment differences computed by the ILP solver (migration plan).}} \\
&  $\mathcal{P}_k$ & \texttt{placement\_before} dictionary& \\
&  $\mathcal{P}_{k+1}$ & \texttt{placement\_after} dictionary& \\

\addlinespace
\textbf{Next State} ($S_{k+1}$)  & $\mathcal{U}_{k+1}$  & \texttt{users\_after} dictionary & \multirow{4}{*}{\pbox{4.2cm}{Reflects the updated system topology and physical allocations after placement execution.}} \\
&  $\mathcal{A}_{k+1}$ & \texttt{apps\_after} dictionary& \\
&  $\mathcal{U}_{k+1}$ & \texttt{users\_after} dictionary& \\
&  $\mathcal{G}_{k+1}$ & \texttt{node\_after}/\texttt{edge\_before} dictionaries& \\
&  $\mathcal{P}_{k+1}$ & \texttt{placement\_after} dictionary& \\


\addlinespace
\textbf{Reward Signal} ($R_k$)  & $f(\mathcal{M}_k,\mathcal{M}_{k+1})$ & \texttt{metrics\_after}/\texttt{metrics\_before} dictionaries & \pbox{4.2cm}{Raw multi-objective telemetry decoupled from reward shaping to ensure training flexibility.} \\
\bottomrule
\end{tabularx}
\end{table*}



\section{Experimental Evaluation and Results}
\label{sect_experimentalevaluationandresults}

This section presents a set of representative configurations for generating datasets that emulate realistic Computing Continuum scenarios. We first introduce a generic baseline scenario built from statistical models widely used in the literature and relying only on the simulator's basic event catalog. We then show how \textsc{ECCOS} can be exploited more effectively by simulating realistic situations that require more complex state modifications, implemented through composed events in the simulator. For each dataset or experiment, we also analyze the temporal evolution of selected metrics in order to illustrate how the simulated scenario changes over time and how the generated traces capture such dynamics.

\begin{table*}[!ht]
\caption{Configurable aspects via statistical distributions in ECCOS, optimal modeling, and bibliographic justification.}
\label{tab:configurable_distributions}
\begin{tabular*}{\textwidth}{@{\extracolsep\fill}p{6.2cm} p{2.2cm} p{2.8cm} p{2.5cm}@{\extracolsep\fill}}
\toprule
\textbf{Configurable Aspect and YAML Path} & \textbf{Optimal Distribution} & \textbf{Realistic Parameters} & \textbf{Bibliographic Reference} \\
\midrule

% INFRASTRUCTURE
Node Capacity (RAM)\newline \texttt{+ infrastructure.node.attributes.ram.distribution} & Lognormal & $\mu=0.1, \sigma=0.5$ & \cite{10.1145/2391229.2391236} \\
\midrule
Edge Capacity (Delay)\newline \texttt{+ infrastructure.edge.attributes.delay.distribution} & Lognormal & $\mu=10.0, \sigma=2.0$ & \cite{10.1145/1879141.1879175} \\
\midrule
Failure Frequency (Nodes / Edges)\newline \texttt{+ infrastructure.node.actions.disable\_node.frequency} & Weibull & $\text{shape}=0.8$ \newline $\text{scale}=300.0$ & \cite{Schroeder_2007} \\
\midrule
Recovery Time (MTTR)\newline \texttt{+ node.actions.disable\_node.impact.distribution\_to\_enable\_node} & Lognormal & $\mu=2.5, \sigma=1.2$ & \cite{Schroeder_2007}\\
\midrule
Node/Network Degradation Magnitude\newline \texttt{+ node.actions.degrade\_node.impact.p\_loss\_dist} & Beta & $\alpha=2.0, \beta=5.0$ & \cite{10.1145/2885497} \\

% APPLICATIONS
\midrule
Microservice Resource Demand (RAM)\newline \texttt{+ app.services.profiles.<profile>.attributes.ram.distribution} & Lognormal & $\mu=2.0, \sigma=1.0$ & \cite{10.1145/3132747.3132772} \\
\midrule
Global Application Popularity\newline \texttt{+ app.popularity.distribution} & Zipf & $a=1.2$ & \cite{749260} \\
\midrule
Popularity Surge Duration\newline \texttt{+ app.actions.surge\_popularity.impact.duration\_dist} & Weibull & $a=1.5$ \newline $\text{scale}=10.0$& \cite{10.1145/1298306.1298309} \\
\midrule
Resource Footprint Variation (Noise)\newline \texttt{+ app.actions.update\_app\_footprint.impact.sigma\_footprint} & Uniform & $\text{low}=0.05$ \newline $\text{high}=0.15$ & \cite{3388242.3388280} \\
\midrule
Geographic Demand Drift\newline \texttt{+ app.actions.geo\_demand\_shift.impact.drift\_dist} & Normal & $\mu=0.0, \sigma=0.05$ & \cite{10.1002/wcm.72} \\

% USERS
\midrule
User Movement Speed\newline \texttt{+ user.mobility.speed} & Uniform & $\text{low}=1.0$ \newline $\text{high}=1.5$ (m/s) & \cite{10.1002/wcm.72} \\
\midrule
Wireless Coverage Radius\newline \texttt{+ user.mobility.coverage\_radius} & Lognormal & $\mu=0.2, \sigma=0.02$ & \cite{Rappaport_2024} \\
\midrule
Initial Request Rate\newline \texttt{+ user.request\_ratio} & Exponential & $\text{scale}=5.0$ & \cite{10.1145/190809.190338} \\
\midrule
Temporary Disconnection Duration\newline \texttt{+ user.actions.suspend\_user.impact.distribution\_to\_resume\_user} & Lognormal & $\mu=1.0, \sigma=0.5$ & \cite{10.1145/2307636.2307658} \\

% GLOBAL SPAWNER
\midrule
New Users Arrival\newline \texttt{+ global\_spawner.actions.new\_user.frequency} & Exponential & $\text{scale}=30.0$ & \cite{10.1145/233008.233034} \\
\midrule
New Apps Arrival\newline \texttt{+ global\_spawner.actions.new\_app.frequency} & Exponential & $\text{scale}=1000.0$ & \cite{10.1145/233008.233034} \\
\bottomrule
\end{tabular*}
\end{table*}


\begin{table*}[!ht]
\caption{Scenario configuration aspects defined in \textsc{ECCOS} by selecting implemented alternatives or by specifying specific parameter values, including bibliographic justification.}
\label{tab:configurable_alternatives}
\begin{tabular*}{\textwidth}{@{\extracolsep\fill}p{6.2cm} p{2.2cm} p{2.8cm} p{2.5cm}@{\extracolsep\fill}}
\toprule
\textbf{Configurable Aspect and YAML Path} & \textbf{Optimal Alternative} & \textbf{Realistic Parameters} & \textbf{Bibliographic Reference} \\
\midrule

% INFRASTRUCTURE
Infrastructure topology \newline \texttt{+ infrastructure.model.name} & \texttt{multi\_tier} & \texttt{cloud.percentatge:2} \newline \texttt{fog.percentatge:10} & \cite{7931566,10.1145/513800.513811,10.1145/316188.316229,587723,10.5555/882459.882563} \\
\midrule
Node and edge attributes model \newline \texttt{+  attributes.<attribute>.mode} & \texttt{centrality\_based} & \texttt{centrality\_type: direct\_proportional} & \cite{PhysRevLett.87.278701} \\
\bottomrule
\end{tabular*}
\end{table*}


\subsection{Dataset Generation}

\subsubsection{Randomized Baseline Scenario}
\subsubsection{Transient Infrastructure Disruption Scenario}
Electric storm with nodes and edges failures
\subsubsection{Localized Crowd-Driven Service Degradation Scenario}
Protest / football match with node and edges degradation and hotspot shifts and popularity changes
\subsubsection{Application Demand Surge Scenario}
Popularity increase and node/edge degradation.

% Esta sección presentará el setup experimental, el hardware utilizado y los parámetros fijos (semillas, tamaños de red).
% El objetivo es demostrar la robustez de ECCOS, no resolver el problema del emplazamiento en sí, sino validar que la herramienta es válida para la comunidad.

\subsection{\sout{Third-party integrations}}

\subsubsection{\sout{Gymnasium}}
\subsubsection{\sout{Enhancing DES Capabilities of the ¿MARIO? Simulator}}











%\subsection{Experimento 1: Escalabilidad y Rendimiento de los Motores de Optimización (Solvers Baseline)}
% OBJETIVO: Demostrar que Ecco-Solver funciona y comparar los baselines.
% DETALLE: Se generarán topologías incrementales (de 50 a 1000 nodos). Para un evento de disrupción masiva constante, se medirá el tiempo de ejecución (overhead) y la calidad de la solución (coste de latencia) del ILP 1, ILP 2, Greedy y Aleatorio.
% QUÉ SE BUSCA: Mostrar que el ILP es óptimo pero inescalable en grandes redes, justificando la necesidad de los datasets de ECCOS para entrenar IAs que logren resultados de ILP en tiempos de Greedy.

%\subsection{Experimento 2: Impacto de las Políticas de Disparo en la Generación de Trazas}
% OBJETIVO: Validar la utilidad de los "Triggers" (solve_all vs solve_on_event_types vs solve_every_t_seconds).
% DETALLE: Se correrá la misma simulación de 10,000 eventos con diferentes políticas. Se medirá cuántas llamadas al solver se ahorran y cómo se acumula la degradación del sistema (QoS drops) cuando el solver se ejecuta con menor frecuencia.
% QUÉ SE BUSCA: Demostrar que ECCOS puede configurar el "paso de tiempo" (step) ideal para modelar problemas de Reinforcement Learning.

%\subsection{Experimento 3: Resiliencia ante Eventos Complejos (Brownouts y Flash Crowds)}
% OBJETIVO: Demostrar el realismo dinámico de Ecco-Generator.
% DETALLE: A lo largo de una línea de tiempo (e.g., 2 horas simuladas), se inyectará una degradación de capacidad en la capa Fog y un aumento súbito de popularidad en una zona geográfica (Flash Crowd).
% QUÉ SE BUSCA: Mediante gráficas temporales (Swimlanes o series de tiempo), ilustrar cómo el sistema entra en sobrecarga y cómo el solver (ej. Greedy) realiza migraciones en cascada para recuperar la estabilidad.

%\subsection{Experimento 4: Caso de Uso - Extracción de Dataset y Topología SFC}
% OBJETIVO: Probar que ECCOS genera la salida prometida.
% DETALLE: Se mostrará un análisis estadístico del dataset en bruto generado (formato, características, balance de clases de eventos).
% QUÉ SE BUSCA: Validar que el encadenamiento de servicios (SFC) se refleja correctamente en el dataset (distribución de microservicios a través del Edge y Cloud).




\section{Conclusions}\label{sec5}



% Resumen de los hitos alcanzados con ECCOS: una herramienta escalable, matemáticamente rigurosa y lista para ser el estándar en generación de datasets en el Continuum.
% Mención a futuras mejoras (e.g., inyección de topologías reales en Ecco-Architect o integración directa con frameworks de PyTorch/TensorFlow).


\sout{Lorem ipsum dolor sit amet, consectetuer adipiscing elit. Ut purus elit, vestibulum ut, placerat ac, adipiscing vitae,
felis. Curabitur dictum gravida mauris. Nam arcu libero, nonummy eget, consectetuer id, vulputate a, magna. Donec
vehicula augue eu neque. Pellentesque habitant morbi tristique senectus et netus et malesuada fames ac turpis egestas.
Mauris ut leo. Cras viverra metus rhoncus sem. Nulla et lectus vestibulum urna fringilla ultrices. Phasellus eu tellus
sit amet tortor gravida placerat. Integer sapien est, iaculis in, pretium quis, viverra ac, nunc. Praesent eget sem vel
leo ultrices bibendum. Aenean faucibus. Morbi dolor nulla, malesuada eu, pulvinar at, mollis ac, nulla. Curabitur
auctor semper nulla. Donec varius orci eget risus. Duis nibh mi, congue eu, accumsan eleifend, sagittis quis, diam.
Duis eget orci sit amet orci dignissim rutrum.

Nam dui ligula, fringilla a, euismod sodales, sollicitudin vel, wisi. Morbi auctor lorem non justo. Nam lacus libero,
pretium at, lobortis vitae, ultricies et, tellus. Donec aliquet, tortor sed accumsan bibendum, erat ligula aliquet magna,
vitae ornare odio metus a mi. Morbi ac orci et nisl hendrerit mollis. Suspendisse ut massa. Cras nec ante. Pellentesque
a nulla. Cum sociis natoque penatibus et magnis dis parturient montes, nascetur ridiculus mus. Aliquam tincidunt
urna. Nulla ullamcorper vestibulum turpis. Pellentesque cursus luctus mauris.}


%\backmatter
\bmsubsection*{Author Contributions}

\textbf{Mireia Jaume:} Data Curation, Formal analysis, Methodology, Software, Validation, Visualization, Writing - review \& editing. \textbf{Isaac Lera:} Conceptualization, Formal analyisis, Funding acquisition, Investigation, Methodology, Project administration, Resources, Supervision, Writing - review \& editing. \textbf{Carlos Guerrero:} Conceptualization, Data curation, Formal analysis, Funding acquisition, Investigation, Methodology, Project administration, Resources, Software, Supervision, Validation, Visualization, Writing - original draft.

\bmsubsection*{Acknowledgments}
The authors acknowledge the use of GPT-5.2 (accessed june-july 2026) to support improve language correction and to correct language typos. All AI-generated suggestions were reviewed, revised, and approved by the authors, who take full responsibility for the accuracy and integrity of the work.

The authors acknowledge the use of Gemini 3.5 Flash and Gemini 3.1 Pro  (accessed may-june 2026) for partial refactoring tasks in the source code of ECCOS and for the detection of errors in the software testing phase.

\bmsubsection*{Financial Disclosure}

This work was supported by Grant PID2024-158637OB-I00, funded by MICIU/AEI/10.13039/501100011033 and by “ERDF A way of making Europe”, “ERDF/EU”.


\bmsubsection*{Supporting Information}

The data generated in this study and the source code of ECCOS are available in a public repository at \url{https://github.com/acsicuib/XXXXXXX/}, under the MIT License.

\bmsubsection*{Conflicts of Interest}

The authors declare no conflicts of interest.

\bibliography{wileyNJD-Chicago}

%\nocite{*}% Show all bib entries - both cited and uncited; comment this line to view only cited bib entries;






\appendix

\bmsection{List of Available Statistical Distributions\label{anex:distributions}}
\vspace*{12pt}


A fundamental design principle of \textsc{ECCO-Generator} is its configuration flexibility. While the framework suggests specific statistical distributions grounded in Computing Continuum and reliability literature to ensure maximum realism, researchers are not locked into these defaults. Every event, its temporality, and its specific behavioral parameters are fully customizable via a YAML configuration file, allowing the substitution of distributions (e.g., swapping a Poisson process for a deterministic interval) without modifying the underlying source code. A complete set of predefined statistical distributions is available to the user (Table~\ref{tab:distributions}).


\begin{table*}[!ht]%
\centering
\caption{Statistical distributions supported by the YAML configuration, inspired by \texttt{numpy.random}.\label{tab:distributions}}
\begingroup
\setlength{\tabcolsep}{2.5pt}
\renewcommand{\arraystretch}{0.72}
\fontsize{6}{6.8}\selectfont
\begin{tabular*}{\textwidth}{@{\extracolsep{\fill}}p{0.22\textwidth}p{0.29\textwidth}p{0.39\textwidth}@{}}
\toprule
\textbf{Distribution} & \textbf{Parameters (Type)} & \textbf{YAML Example} \\
\midrule
Constant (Numeric Literal) & None; direct constant value & \begin{tabular}[t]{@{}l@{}} \texttt{distribution:~100} \end{tabular} \\
\midrule
Normal (Gaussian) & \begin{tabular}[t]{@{}l@{}} \texttt{loc} (float) \\ \texttt{scale} (float) \end{tabular} & \begin{tabular}[t]{@{}l@{}} \texttt{distribution:} \\ \texttt{~~type:~normal} \\ \texttt{~~loc:~10.0} \\ \texttt{~~scale:~2.5} \end{tabular} \\
\midrule
Uniform & \begin{tabular}[t]{@{}l@{}} \texttt{low} (float) \\ \texttt{high} (float) \end{tabular} & \begin{tabular}[t]{@{}l@{}} \texttt{distribution:} \\ \texttt{~~type:~uniform} \\ \texttt{~~low:~0.0} \\ \texttt{~~high:~1.0} \end{tabular} \\
\midrule
Exponential & \texttt{scale} (float) & \begin{tabular}[t]{@{}l@{}} \texttt{distribution:} \\ \texttt{~~type:~exponential} \\ \texttt{~~scale:~5.0} \end{tabular} \\
\midrule
Poisson & \texttt{lam} (float) & \begin{tabular}[t]{@{}l@{}} \texttt{distribution:} \\ \texttt{~~type:~poisson} \\ \texttt{~~lam:~4.0} \end{tabular} \\
\midrule
Pareto & \begin{tabular}[t]{@{}l@{}} \texttt{a} (float) \\ \texttt{size} (string/expression) \end{tabular} & \begin{tabular}[t]{@{}l@{}} \texttt{distribution:} \\ \texttt{~~type:~pareto} \\ \texttt{~~a:~1.5} \\ \texttt{~~size:~"\{num\_nodes\}"} \end{tabular} \\
\midrule
Integers & \begin{tabular}[t]{@{}l@{}} \texttt{low} (integer) \\ \texttt{high} (integer) \end{tabular} & \begin{tabular}[t]{@{}l@{}} \texttt{distribution:} \\ \texttt{~~type:~integers} \\ \texttt{~~low:~1} \\ \texttt{~~high:~10} \end{tabular} \\
\midrule
Choice (Set Selection) & \texttt{a} (string/expression) & \begin{tabular}[t]{@{}l@{}} \texttt{distribution:} \\ \texttt{~~type:~choice} \\ \texttt{~~a:~"\{active\_edges\}"} \end{tabular} \\
\midrule
Binomial & \begin{tabular}[t]{@{}l@{}} \texttt{n} (integer) \\ \texttt{p} (float) \end{tabular} & \begin{tabular}[t]{@{}l@{}} \texttt{distribution:} \\ \texttt{~~type:~binomial} \\ \texttt{~~n:~10} \\ \texttt{~~p:~0.5} \end{tabular} \\
\midrule
Gamma & \begin{tabular}[t]{@{}l@{}} \texttt{shape} (float) \\ \texttt{scale} (float) \end{tabular} & \begin{tabular}[t]{@{}l@{}} \texttt{distribution:} \\ \texttt{~~type:~gamma} \\ \texttt{~~shape:~2.0} \\ \texttt{~~scale:~1.0} \end{tabular} \\
\midrule
Beta & \begin{tabular}[t]{@{}l@{}} \texttt{a} (float) \\ \texttt{b} (float) \end{tabular} & \begin{tabular}[t]{@{}l@{}} \texttt{distribution:} \\ \texttt{~~type:~beta} \\ \texttt{~~a:~0.5} \\ \texttt{~~b:~0.5} \end{tabular} \\
\midrule
Lognormal & \begin{tabular}[t]{@{}l@{}} \texttt{mean} (float) \\ \texttt{sigma} (float) \end{tabular} & \begin{tabular}[t]{@{}l@{}} \texttt{distribution:} \\ \texttt{~~type:~lognormal} \\ \texttt{~~mean:~0.0} \\ \texttt{~~sigma:~1.0} \end{tabular} \\
\midrule
Weibull & \texttt{a} (float) & \begin{tabular}[t]{@{}l@{}} \texttt{distribution:} \\ \texttt{~~type:~weibull} \\ \texttt{~~a:~1.5} \\ \texttt{~~scale:~10.0} \end{tabular} \\
\midrule
Geometric & \texttt{p} (float) & \begin{tabular}[t]{@{}l@{}} \texttt{distribution:} \\ \texttt{~~type:~geometric} \\ \texttt{~~p:~0.35} \end{tabular} \\
\midrule
Dirichlet & \texttt{alpha} (list of floats) & \begin{tabular}[t]{@{}l@{}} \texttt{distribution:} \\ \texttt{~~type:~dirichlet} \\ \texttt{~~alpha:~[10.0,~5.0,~2.0]} \end{tabular} \\
\midrule
Multinomial & \begin{tabular}[t]{@{}l@{}} \texttt{n} (integer) \\ \texttt{pvals} (list of floats) \end{tabular} & \begin{tabular}[t]{@{}l@{}} \texttt{distribution:} \\ \texttt{~~type:~multinomial} \\ \texttt{~~n:~20} \\ \texttt{~~pvals:~[1.0,~1.0,~1.0]} \end{tabular} \\
\midrule
Triangular & \begin{tabular}[t]{@{}l@{}} \texttt{left} (float) \\ \texttt{mode} (float) \\ \texttt{right} (float) \end{tabular} & \begin{tabular}[t]{@{}l@{}} \texttt{distribution:} \\ \texttt{~~type:~triangular} \\ \texttt{~~left:~0.0} \\ \texttt{~~mode:~0.5} \\ \texttt{~~right:~1.0} \end{tabular} \\
\midrule
Rayleigh & \texttt{scale} (float) & \begin{tabular}[t]{@{}l@{}} \texttt{distribution:} \\ \texttt{~~type:~rayleigh} \\ \texttt{~~scale:~1.0} \end{tabular} \\
\midrule
Standard Cauchy & None & \begin{tabular}[t]{@{}l@{}} \texttt{distribution:} \\ \texttt{~~type:~standard\_cauchy} \end{tabular} \\
\midrule
Random Float [0.0, 1.0) & None & \begin{tabular}[t]{@{}l@{}} \texttt{distribution:} \\ \texttt{~~type:~random} \end{tabular} \\
\bottomrule
\end{tabular*}
\endgroup
\end{table*}


\bmsection{YAML Configuration File Schema for Scenario Definition\label{app1.1a}}

This appendix describes the exact format of the YAML configuration file used to define an \textsc{ECCOS} simulation scenario. This file specifies the scenario structure and stochastic behavior consumed by the simulator, including the infrastructure, applications, users, global spawners, and event definitions. The configuration filename is defined in the simulator's \texttt{constants.py} file, although it can also be provided explicitly through the command line when launching a simulation. By default, \textsc{ECCOS} expects this file to be named \texttt{scenario\_config.yaml}.

To provide fine-grained control over experimental reproducibility, the scenario configuration YAML file supports a dedicated \texttt{seeds} block. First, users may define either a single master seed, which automatically derives deterministic, sequential sub-seeds across all simulation domains. Alternativelly, users may explicitly specify independent integer seeds for individual conceptual domains (\texttt{graph}, \texttt{app}, \texttt{user}, and \texttt{event}). When explicit domain seeds are provided, they override the master derivation for that specific scope.

\begin{lstlisting}[caption={YAML configuration for stochastic experimental reproductibility through seeds definition.},label=lst:user_event, basicstyle=\fontsize{8}{10}\selectfont\ttfamily]
seeds:
  master: 42
  graph: 42
  app: 43
  user: 44
  event: 45
\end{lstlisting}



\bmsubsection{Infrastructure Definition}
\label{anex:infrastructuredefinition}
The infrastructure definition in Ecco-Architect is driven by the number of nodes (\texttt{num\_nodes}), the topological model for the generation of the graph structure (\texttt{model}), and the mapping strategy to map graph nodes to the computing continuum layers (\texttt{model.layer.mode}): cloud, fog and edge. The topological model and the mapping of layers requires of additional parameters in some cases. 

\begin{lstlisting}[caption={YAML configuration for infrastructure definition.},label=lst:user_event, basicstyle=\fontsize{8}{10}\selectfont\ttfamily]
infrastructure:
  num_nodes: 50
  model:
    name: scale_free                   # Values (erdos_renyi | scale_free | spatial | fat_tree | 
                                       # tree | multi_tier)
    m: 2                               # For scale_free. Number of edges to attach from a new node
                                       # to existing nodes
    radius: 0.1                        # For spatial. Distance radius for connecting nodes in the
                                       # geometric graph
    branching_r:  2                    # For tree/fat_tree. Branching factor of the tree (number of
                                       # children per node)
    branching_h: 4                     # For tree/fat_tree. Height or depth of the tree
    cloud:
      percentage: 2                    # Percentage of total nodes assigned to the cloud layer
    fog:
      percentage: 10                   # Percentage of total nodes assigned to the fog layer
      m: 2                             # Number of edges to attach per new fog node in its scale-
                                       # free topology (default: 2)
    edge:
      radius: 0.1                      # Connection distance radius for the spatial edge layer 
                                       # (default: 0.1)
    layer:
      mode: centrality_based_layer     # Values (centrality_based_layer | depth_based_layer)
      thresholds:
        cloud_min: 0.1                 # For centrality_based_layer Minimum centrality score (0.0 
                                       # to 1.0) to be considered 'cloud'
        fog_min: 0.02                  # For centrality_based_layer Minimum centrality score to be 
                                       # considered 'fog' (nodes below this become 'edge')
        cloud_max: 0                   # For depth_based_layer.  Only applies to 'tree'/'fat_tree' 
                                       # models. Maximum depth level to be considered 'cloud' 
                                       # (usually 0 for root)
        fog_max: 2                     # For depth_based_layer. Only applies to 'tree'/'fat_tree' 
                                       # models. Maximum depth level to be considered 'fog' (nodes 
                                       # deeper than this become 'edge')
\end{lstlisting}


\textsc{ECCOS} characterizes nodes and edges with unlimited user-defined attributes (\texttt{node.attributtes}, \texttt{edge.attributtes}) to describe resources, operational properties\ldots The configuration must define how its initial values are assigned across the generated topology (\texttt{node.attributtes.<name>.mode}) and general bounds values (\texttt{node.attributtes.<name>.min/max}). Additionally, the attribute distribution allows to define the statistical model to generate the values. In this block, specific layer definitios can overwrite the general definition, allowing to express a specific configuration for the nodes/edges of a given layer (\texttt{node.attributtes.<name>.distribution.cloud/fog/edge}). 



\begin{lstlisting}[caption={YAML configuration for the attributes definition of the infrastructure entities.},label=lst:user_event, basicstyle=\fontsize{8}{10}\selectfont\ttfamily]
infrastructure:
  node:
    attributes:
      cpu:
        mode: centrality_based                  # Values (uniform | centrality_based | layered)
        centrality_type: direct_proportional    # For centrality_based. Values (direct_proportional
                                                # | inverted proportional)
        min: 4.0
        max: 16.0
        distribution:
          type: pareto
          a: 1.16
          cloud:                                # Overrites distribution and max for cloud nodes.
            type: uniform
            low: 128.0
            high: 512.0
            max: 1024.0    
      ram:
        mode: layered
        min: 1.0                                # General bounds
        max: 1024.0         
        distribution:
          cloud:
            type: uniform
            low: 256.0
            high: 512.0
            min: 200.0                          # Overwrites values
            max: 1024.0
          fog:
            type: normal
            loc: 32.0
            scale: 8.0
            min: 8.0                            # <-- Overwrites values
            max: 64.0
          edge:                                 # No explicit min/max values, it inherits the 
                                                # global 1.0 and 1024.0 bounds
            type: uniform
            low: 2.0
            high: 8.0
            
\end{lstlisting}

\bmsubsection{Infrastructure Events}
\label{anex_infrastructureevents}

Events to disable nodes and edges (\texttt{disable\_node}, \texttt{disable\_edge}) are configured in terms of their frequency and their impact. Frequency is modeled through the use of a statistical distribution to define their inter-arrival time. The impact of the disable events is measured in terms of the MTTR. Thus, both events hast the block \texttt{distribution\_to\_enable\_node} in their impact block, also modeled with a statistical distribution.

\begin{lstlisting}[caption={YAML configuration for disable events of infrastructure elements},label=lst:global_spawner_yaml, basicstyle=\fontsize{8}{10}\selectfont\ttfamily]
infrastructure:
  node:
    actions:
      disable_node:
        frequency:
          type: exponential
          scale: 300
        impact:
          distribution_to_enable_node:
            type: uniform
            low: 0
            high: 100
  edge:
    actions:
      disable_edge:
        frequency:
          type: exponential
          scale: 300
        impact:
          distribution_to_enable_edge:
            type: uniform
            low: 0
            high: 100            
\end{lstlisting}


Degradation events such as \texttt{degrade\_node} and  \texttt{congest\_edge} are also modeled through a statistical distribution for their inter-arrival times (\texttt{frequency}). But in this case, the impact is not only configured with the MTTR distribution (\texttt{impact.distribution\_to\_restore\_node/distribution\_to\_clear\_edge}), but also by the list of attributes that are degraded or multiplied (\texttt{impact.target\_attributes/target\_attributes\_loss/target\_attributtes\_multiplier}), and the statistical distribution that govern those losses/multipliers (\texttt{impact.p\_loss\_dist/multiplier\_dist}). 


\begin{lstlisting}[caption={YAML configuration for degrade events of infrastructure elements},label=lst:global_spawner_yaml, basicstyle=\fontsize{8}{10}\selectfont\ttfamily]
infrastructure:
  node:
    actions:
      degrade_node:
        frequency:
          type: exponential
          scale: 150
        impact:
          target_attributes: ["ram", "cpu"]
          p_loss_dist:
            type: beta
            a: 2
            b: 5
          distribution_to_restore_node:
            type: exponential
            scale: 50
  edge:
    actions:
      congest_edge:
        frequency:
          type: exponential
          scale: 150
        impact:
          target_attributes_loss: ["bandwidth"]
          target_attributes_multiplier: ["delay"]
          p_loss_dist:
            type: beta
            a: 2
            b: 5
          multiplier_dist:
            type: lognormal
            mean: 0.5
            sigma: 0.2
          distribution_to_clear_edge:
            type: exponential
            scale: 50        
\end{lstlisting}



\bmsubsection{User Definition}
\label{anex:userdefinition}

Users are modelled by the initial number of users (\texttt{num\_users}), the statistical distribution to determine their request ratios (\texttt{request\_ratio}), the spatial size of the simulated area where the users can move arround (\texttt{spatial\_region.width/height}), the streets of this simulated area (\texttt{spatial\_region.num\_vertical\_streets/num\_horizontal\_streets}), and the model for the initial position of the users (\texttt{spatial\_distribution}).


\begin{lstlisting}[caption={YAML configuration for users creation.},label=lst:user_event, basicstyle=\fontsize{8}{10}\selectfont\ttfamily]
user:
  num_users: 2
  request_ratio:
    type: exponential
    scale: 5
  spatial_distribution:
    model: thomas_cluster          # Values (random_uniform | thomas_cluster)
    hotspots: 5                    # only for thomas_cluster
    spread: 0.05                   # only for thomas_cluster
  spatial_region:
    width: 2.0
    height: 2.0
    num_vertical_streets: 10
    num_horizontal_streets: 10
\end{lstlisting}

Additionally, user mobility patterns are configured through the mobility model(\texttt{model}), the direction-change or turning probabilities (\texttt{turn\_probabilities}), the movement speed (\texttt{speed}), and the distance that defines the coverage range to the connection node for each user (\texttt{coverage\_radius}).


\begin{lstlisting}[caption={YAML configuration for the definition of the users mobility mode.},label=lst:user_event, basicstyle=\fontsize{8}{10}\selectfont\ttfamily]
user:
  mobility:
    model: random              # Values (static | random | manhattan)
    turn_probabilities:
      straight: 0.33
      left: 0.33               # Only for manhattan
      right: 0.34              # Only for manhattan
    speed:
      type: uniform
      low: 0.005
      high: 0.02
    coverage_radius:
      type: normal
      mean: 0.2
      sigma: 0.02
\end{lstlisting}  





\bmsubsection{User Events}
\label{anex_userevents}

The user removal event (\texttt{remove\_user}) is modeled through its frequency and the parametrization of the corresponding statistical distribution, and therefore does not require an impact configuration. Similarly, user mobility across the two-dimensional simulation plane is modeled by the user movement event (\texttt{move\_user}), together with its corresponding frequency distribution. Its impact is determined by the mobility configuration established during user definition and creation.


\begin{lstlisting}[caption={YAML configuration for basic user events.},label=lst:user_event, basicstyle=\fontsize{8}{10}\selectfont\ttfamily]
user:
  actions:
    remove_user:
      distribution:
        type: exponential
        scale: 300
    move_user:
      frequency:
        type: exponential
        scale: 120
\end{lstlisting}


Temporary user disconnections are triggered by the suspend-user event (\texttt{suspend\_user}). In addition, an automatic reconnection event is inserted into the event queue to simulate device recovery. Both the frequency of user disconnections and the time required to restore connectivity (\texttt{distribution\_to\_resume\_user}) are modeled through their corresponding statistical distributions.

\begin{lstlisting}[caption={YAML configuration for user disconection event.},label=lst:user_event, basicstyle=\fontsize{8}{10}\selectfont\ttfamily]
user:
  actions:
    suspend_user:
      frequency:
        type: exponential
        scale: 150
      impact:
        distribution_to_resume_user:
          type: lognormal
          mean: 1.0
          sigma: 0.5
\end{lstlisting}

Temporary changes in users' request ratios are modeled through the \texttt{change\_request\_ratio} event and its inter-arrival frequency distribution. The event impact is configured through the probability of producing either a burst, which increases requests, or a valley, which decreases them (\texttt{increase\_probability}), together with the multiplier magnitude of the request-ratio modification (\texttt{multiplier}) and the time required to restore the initial ratio (\texttt{restore\_ratio}).

\begin{lstlisting}[caption={YAML configuration for activity changes user event.},label=lst:user_event, basicstyle=\fontsize{8}{10}\selectfont\ttfamily]
user:
  actions:
    change_request_ratio:
      frequency:
        type: exponential
        scale: 90
      impact:
        increase_probability: 0.5
        multiplier:
          type: uniform
          low: 1.5
          high: 3.0
        restore_ratio:
          type: lognormal
          mean: 1.0
          sigma: 0.5
\end{lstlisting}

\bmsubsection{Application Definition}
\label{anex:applicationdefinition}


The simulator must know the number of applications to generate at the beginning of the simulation (\texttt{num\_apps}). At the same time, the configuration defines the range of services that may compose each application (\texttt{services.num\_services\_range}) and how these applications are handled during placement (\texttt{architecture}): either as a single monolithic block that must be deployed on one node (\texttt{monolithic}), or as separable microservices that can be placed across different nodes (\texttt{microservices}). For microservice-based applications, the service composition can be configured (\texttt{topology\_model}) either as a Service Function Chain (SFC), which imposes a strict sequential order among services, or as a Directed Acyclic Graph (DAG). In the latter case, the configuration must specify the graph-generation parameters, including the initial number of nodes (\texttt{initial\_nodes})  and the number of edges attached per node (\texttt{edges\_per\_node}).

\begin{lstlisting}[caption={YAML configuration for application creation.},label=lst:user_event, basicstyle=\fontsize{8}{10}\selectfont\ttfamily]
app:
  num_apps: 20
  architecture: monolithic                  # Values (monolithic | microservice)
  services:
    num_services_range: [2, 6]
    topology_model: directed_scale_free     # Values (sfc | directed_scale_free)
    scale_free_settings:                    # Only for directed_scale_free
      initial_nodes: 2
      edges_per_node: 2
\end{lstlisting}

The simulator defines the network resource consumptions as attributtes associated to the application (\texttt{app.services.network.<attribute\_name>}). On contrary, the computational resource consumptions are associated to services through the use of profile definitios (\texttt{app.services.profiles.<profile\_name>.<attribute\_name>}). In the later case, the profiles are selected uniformerly random through the assigned probability (\texttt{prob}).

\begin{lstlisting}[caption={YAML configuration for resource consumptions and profiles definition of the services.},label=lst:user_event, basicstyle=\fontsize{8}{10}\selectfont\ttfamily]
app:
  services:
    network:
      delay:
        distribution:
          type: uniform
          low: 10.0
          high: 50.0
      bandwidth:
        distribution:
          type: uniform
          low: 10.0
          high: 100.0
    profiles:
      iot:
        prob: 0.6
        attributes:
          cpu:
            distribution:
              type: lognormal
              mean: 0.1
              sigma: 0.5
          ram:
            distribution:
              type: lognormal
              mean: 0.1
              sigma: 0.5
      video:
        prob: 0.4
        attributes:
          cpu:
            distribution:
              type: lognormal
              mean: 2.0
              sigma: 1.0
          ram:
            distribution:
              type: lognormal
              mean: 3.0
              sigma: 1.5
\end{lstlisting}


Finally, applications must define their popularity (\texttt{popularity}) through a statistical distribution. The popularity is the probability with which applications are requested by newly created users or by the initial users at the begining of the simulation. The configuration also allows a subset of applications to be marked as local (\texttt{local\_app\_ratio}), increasing their popularity within their corresponding influence radius (\texttt{local\_radius\_influence}). To avoid creating applications without immediate demand, the simulator also defines the number of new users to be generated and associated with each newly created application (\texttt{num\_new\_users}).

\begin{lstlisting}[caption={YAML configuration for application popularity lifecycle.},label=lst:user_event, basicstyle=\fontsize{8}{10}\selectfont\ttfamily]
app:
  popularity:
    distribution:
      type: zipf
      a: 1.2
    local_app_ratio: 0.3
    local_radius_influence: 0.15
  num_new_users:
    type: integers
    low: 1
    high: 6
\end{lstlisting}



\bmsubsection{Application Events}
\label{anex_appsevents}

Application deletion event (\texttt{remove\_app}) is simply modelled with its frequency distribution. 

\begin{lstlisting}[caption={YAML configuration for removal application event.},label=lst:user_event, basicstyle=\fontsize{8}{10}\selectfont\ttfamily]
app:
  actions:
    remove_app:
      frequency:
        type: uniform
        low: 0
        high: 300
\end{lstlisting}

The events to change the application workloads (\texttt{update\_app\_footprint}) and network usages (\texttt{update\_app\_network}) are identically modeled through their frequency, and the impact based on the distribution for the stadard deviation of the applied noise (\texttt{sigma\_footprint}).

\begin{lstlisting}[caption={YAML configuration for application workload changes events.},label=lst:user_event, basicstyle=\fontsize{8}{10}\selectfont\ttfamily]
app:
  actions:
    update_app_footprint:
      frequency:
        type: exponential
        scale: 100
      impact:
        sigma_footprint:
          type: uniform
          low: 0.05
          high: 0.15
    update_app_network:
      frequency:
        type: exponential
        scale: 150
      impact:
        sigma_net:
          type: uniform
          low: 0.05
          high: 0.15
\end{lstlisting}



The event that changes the topological structure of applications in microservice-defined architectures is \texttt{update\_app\_topology}. As in the other cases, this event is modeled by the individual frequency of each entity and by its impact, which is measured in terms of the minimum and maximum number of services that remain after the application update (\texttt{l\_min} and \texttt{l\_max\_global}).

\begin{lstlisting}[caption={YAML configuration for application lifecycle event.},label=lst:user_event, basicstyle=\fontsize{8}{10}\selectfont\ttfamily]
app:
  actions:
    update_app_topology:
      frequency:
        type: exponential
        scale: 120
      impact:
        l_min:
          type: integers
          low: 2
          high: 3
        l_max_global:
          type: integers
          low: 6
          high: 7
\end{lstlisting}


Changes in application popularity, restricted to newly created users, are defined through the \texttt{surge\_popularity} and \texttt{drop\_popularity} events, which respectively increase or decrease an application's popularity. This popularity change can be transient or permanent, as modeled by the statistical distribution in \texttt{transient\_prob}; for transient changes, the duration of the updated popularity is defined by \texttt{duration\_dist}.


\begin{lstlisting}[caption={YAML configuration for application popularity events.},label=lst:user_event, basicstyle=\fontsize{8}{10}\selectfont\ttfamily]
app:
  actions:
    surge_popularity:           # or drop_popularity:
      frequency:
        type: exponential
        scale: 200
      impact:
        transient_prob:
          type: uniform
          low: 0.7
          high: 0.7
        duration_dist:
          type: weibull
          a: 1.5
          scale: 10.0
\end{lstlisting}


Changes in the location of the hot spot for localized applications can be triggered through the \texttt{geo\_demand\_shift} event, which moves the hot spot either by a random jump or by an incremental movement, according to the \texttt{jump\_prob} parameter. In the case of drift-based movement of the hot spot, the statistical distribution in \texttt{drift\_dist} determines the shift from the current coordinates.


\begin{lstlisting}[caption={YAML configuration for hot spot changes for localized applications.},label=lst:user_event, basicstyle=\fontsize{8}{10}\selectfont\ttfamily]
app:
  actions:
    geo_demand_shift:
      frequency:
        type: exponential
        scale: 50
      impact:
        jump_prob:
          type: uniform
          low: 0.2
          high: 0.2
        drift_dist:
          type: normal
          mean: 0.0
          sigma: 0.05
\end{lstlisting}





\bmsubsection{Global Spawner Events}
\label{anex_globalevents}

Creating new users and applications are global simulation events that represents the external arrival of a new entities into the scenario. Therefore, these events are defined in the global configuration entries \texttt{global\_spawner.actions.new\_user} and \texttt{global\_spawner.actions.new\_app}, rather than inside the respective action blocks: \texttt{user.actions} or \texttt{app.actions}.

Unlike domain-specific events where the target entity is implicit, the global orchestration block \texttt{global\_spawner} requires \texttt{type\_object} to dynamically route events to their appropriate domain managers (the UserSet or ApplicationSet catalogs). Simultaneously, \texttt{action\_type} maps the configuration event directly to the specific internal handler method via runtime reflection; this separates declarative configuration naming from code implementation. This separation for the specific case of global orchestration enables complex, cross-domain user-defined composite workflows.

The \texttt{type\_object} attribute also supports the special value \texttt{type\_object:global} for user-defined composed events (Section~\ref{anex:usercomposedevents}).


\begin{lstlisting}[caption={YAML configuration for global entity creation events},label=lst:global_spawner_yaml, basicstyle=\fontsize{8}{10}\selectfont\ttfamily]
global_spawner:
  actions:
    new_user:
      type_object: user             # Values: (user | app | global)
      action_type: new_user         # Values: (new_user | new_app)
      frequency:
        type: exponential
        scale: 30
    new_app:
      type_object: app
      action_type: new_app
      frequency:
        type: exponential
        scale: 3600
\end{lstlisting}





\bmsubsection{User Composition of Encapsulated Events}
\label{anex:usercomposedevents}

An encapsulated event is defined, like other events, as an element of the \texttt{actions} attribute, using a user-defined name (\texttt{actions.<event\_name>}) and its corresponding inter-arrival frequency (\texttt{frequency}). In addition, its impact specification includes a list of events under the \texttt{impact.composed\_of} attribute. Because of the generic nature of the composition model, each event in this list must explicitly declare the \texttt{type\_object} and \texttt{action\_type} attributes. The \texttt{type\_object} attribute indicates the entity type that handles the event, whereas \texttt{action\_type} specifies the exact base action to trigger. 

Table~\ref{tab:composite_actions} includes an exhaustive enumeration of all the \texttt{action\_type} alternatives. It is important to note that this table includes more actions than the set presented as base configurable actions in the main text. This is intentional: some entries correspond to derived actions that are internally triggered by other base events, but may still be useful when defining composed events. For example, \texttt{remove\_user\_by\_requested\_app} is derived from \texttt{remove\_app}, since removing an application also requires removing all users requesting it. Similarly, \texttt{revive\_edge} is automatically triggered after \texttt{disable\_edge} according to the configured recovery process. These derived actions are typically less useful as standalone base-event configurations, but they provide additional flexibility when constructing user-defined composed actions.


\begin{center}
\begin{table*}[!ht]%
\caption{Taxonomy of composite action types (\texttt{action\_type}) categorized by simulation entity domain (\texttt{type\_object}).\label{tab:composite_actions}}
\begin{tabular*}{\textwidth}{@{\extracolsep\fill}ll p{6.0cm} p{5.0cm}@{\extracolsep\fill}}
\toprule
\textbf{Domain} & \textbf{Action}  & \textbf{Primary Simulation Impact}  & \textbf{Operational Context}   \\
\textbf{(\texttt{type\_object})} & \textbf{(\texttt{action\_type})}  &   &    \\
\midrule
\multirow{6}{*}{\textbf{\texttt{user}}} 
 & \texttt{move\_user} & Displaces $(X, Y)$ coordinates based on mobility models. & Mobile user tracking, handover tests. \\
 & \texttt{change\_request\_ratio} & Modifies workload generation frequency multiplier. & Flash crowds, local demand spikes. \\
 & \texttt{suspend\_user} & Temporarily halts user session and traffic generation. & Network outages, disconnections. \\
 & \texttt{new\_user} & Instantiates or reactivates a user entity dynamically. & Diurnal workload scaling, arrivals. \\
 & \texttt{remove\_user} & Permanently deletes user instance and pending events. & Session termination, cleanup. \\
 & \texttt{remove\_user\_by\_requested\_app} & Batch removal of users linked to a specific application. & Cascading app decommissioning. \\
\midrule
\multirow{7}{*}{\textbf{\texttt{app}}} 
 & \texttt{surge\_popularity} & Triggers transient upward surge in request volume. & Viral campaigns, sudden traffic peaks. \\
 & \texttt{drop\_popularity} & Induces transient decrease in request popularity. & Off-peak transitions, user decay. \\
 & \texttt{geo\_demand\_shift} & Relocates app geographic focus via jump or drift. & Regional demand migration. \\
 & \texttt{update\_app\_footprint} & Modifies service CPU, RAM, and storage requirements. & Vertical scaling, heavy processing. \\
 & \texttt{update\_app\_network} & Alters bandwidth requirements and latency constraints. & QoS policy updates, SLA adjustments. \\
 & \texttt{update\_app\_topology} & Mutates service graph within $[l_{\min}, l_{\max}]$ bounds. & Horizontal architectural adaptation. \\
 & \texttt{new\_app} / \texttt{remove\_app} & Dynamically spawns or destroys application graphs. & Lifecycle orchestration, multi-tenancy. \\
\midrule
\multirow{3}{*}{\textbf{\texttt{graph\_node}}} 
 & \texttt{disable\_node} & Disconnects server, triggering service migrations. & Hardware failures, power outages. \\
 & \texttt{degrade\_node} & Reduces available server CPU/RAM by loss percentage. & Resource exhaustion, overload. \\
 & \texttt{revive\_node} / \texttt{restore\_node} & Restores server to full operational capacity. & Maintenance completion, recovery. \\
\midrule
\multirow{3}{*}{\textbf{\texttt{graph\_edge}}} 
 & \texttt{disable\_edge} & Severs link, forcing routing and shortest-path recalculation. & Fiber cuts, wireless link loss. \\
 & \texttt{congest\_edge} & Degrades link bandwidth and inflates packet loss/latency. & Network congestion, DDoS attacks. \\
 & \texttt{revive\_edge} / \texttt{clear\_edge} & Restores network link to standard operating parameters. & Link restoration, congestion clearing. \\
\bottomrule
\end{tabular*}

\end{table*}
\end{center}

\begin{lstlisting}[escapechar=\%,caption={YAML configuration of a combined move-increase activity user event.},label=lst:yaml_catastrophic, basicstyle=\fontsize{8}{10}\selectfont\ttfamily]
user:
  actions:
    move_and_increse_request:
      frequency:
        type: exponential
        scale: 600
      impact:
        composed_of:                            # The user moves and [...]
          - action_type: move_user              # Values listed in Table %\ref{tab:composite_actions}%
            type_object: user                   # Values (user | app | graph_node | graph_edge)
            impact_params:
              target_resolution:
                mode: self
          - action_type: change_request_ratio   # [...] the user increase request ratio
            type_object: user
            impact_params:
              target_resolution:
                mode: self    
              multiplier:
                type: uniform
                low: 1.5
                high: 3.0
\end{lstlisting}

Each event listed in \texttt{composed\_of} must also model its own impact, in the same way as when that event is configured independently. In encapsulated events, this impact configuration is placed under the \texttt{impact\_params} attribute of each composed sub-action. The expected subattributes are therefore the same as those required by the corresponding standalone event; only their location in the YAML hierarchy changes. Additionally, the new parameter \texttt{target\_resolution} is used to select the entities to apply the event.

This \texttt{target\_resolution} attribute supports the following modes: (a) \texttt{self}, the default behavior. The sub-action is applied strictly to the same entity instance that generated the macro-event; (b) \texttt{random}, a target is selected following a uniform random distribution from the pool of all active instances of the specified \texttt{type\_object}; (c) \texttt{id}, the sub-action is statically bound to a specific instance identifier (e.g., \texttt{"Node\_24"}); (d) \texttt{intelligent}, the target resolution is delegated to the action handler itself, which leverages its context to infer the target (e.g., a \texttt{disable\_node} action receiving an Application ID will dynamically discover and disable the specific node hosting that application); (e) \texttt{group}, which allows a sub-action to affect a cluster of $N$ elements simultaneously to model a wide-spread phenomena.

The \texttt{group\_config} block is required for granular configuration of the \texttt{target\_resolution: group} mode. The size of the cluster (\texttt{num\_elements}) can be defined either as an absolute integer or as a percentage using a floating-point value. The selection strategy of the $N$ elements in the group is driven by the \texttt{group\_config.strategy} directive, which supports explicit: (a) \texttt{list} selection, (b) \texttt{random} sampling, (c) \texttt{self\_random} hybrid approach (the originating instance plus $N-1$ random instances), or the contextual proximity model, (d) \texttt{self\_proximity} and (e) \texttt{random\_proximity}, to model localized phenomena.

Note that assigning \texttt{type\_object: global} is strictly valid only at the level of parent macro-events defined in the \texttt{global\_spawner}. In this context, the global domain acts as a central scheduling timer whose sole architectural purpose is to trigger composite workflows (\texttt{composed\_of}). Conversely, within the \texttt{composed\_of} block itself, assigning \texttt{type\_object: global} is architecturally invalid; every sub-action must explicitly delegate execution to a specific domain manager (\texttt{user}, \texttt{app}, \texttt{graph\_node}, or \texttt{graph\_edge}) where the mathematical and operational action handlers reside. Furthermore, when configuring ex novo instantiation actions (such as \texttt{new\_user} or \texttt{new\_app}), the target resolution mechanism (\texttt{target\_resolution}) must be omitted entirely, as these actions dynamically spawn new entities from scratch without referencing any pre-existing target identifiers.

The following YAML snippet illustrates the versatility of the proposed framework. Listing~\ref{lst:yaml_catastrophic} demonstrates a Catastrophic Failure macro-event. When triggered, it intelligently disables the node related to the origin entity, and concurrently creates a popularity surge in three logically proximate applications to simulate the reallocation of user traffic.


\begin{lstlisting}[escapechar=\%,caption={YAML configuration of a catastrophic failure encapsulated event},label=lst:yaml_catastrophic, basicstyle=\fontsize{8}{10}\selectfont\ttfamily]
global_spawner:
  actions:
    catastrophic_failure_and_surge:
      type_object: global                        # Values: (user | app | global)
      action_type: catastrophic_failure          # Only functional required when user or app type 
                                                 # object, in other case is only informative or 
                                                 # descriptive.
      frequency:
        type: exponential
        scale: 120
      impact:
        composed_of:                            # 1. Cross-entity intelligent targeting
          - type_object: graph_node             # Values (user | app | graph_node | graph_edge)
            action_type: disable_node           # Values listed in Table %\ref{tab:composite_actions}%
            impact_params:
              target_resolution:
                mode: intelligent               # Values (self | random | id | intelligent | group)
#               id: "Node_67"                   # Only for mode:id
              distribution_to_enable_node: 
                type: uniform
                low: 10
                high: 30
          - type_object: app                    # 2. Group targeting based on logical proximity
            action_type: surge_popularity
            impact_params:
              target_resolution:
                mode: group                     # Values (self | random | id | intelligent | group)
                group_config:                   # Only for mode:group
                  strategy: random_proximity    # Values (random | list | self_random | self_proximity
                                                # | random_proximity) 
                  num_elements: 3
#                 list:                         # Only for strategy:list
#                   - "App_0"
#                   - "App_2"
#                   - "App_5"
              transient_prob: 
                type: uniform
                low: 0.9
                high: 0.9
\end{lstlisting}


The generic nature of the composed-action definition also enables users to define multiple variants of the same event, each governed by different statistical behaviors. In this way, a single logical phenomenon can be represented through several stochastic profiles without modifying the simulator implementation. For example, Listing~\ref{lst:yaml_tworegimes} configures two stochastic regimes for node degradation by defining two encapsulated events that both trigger the same \texttt{degrade\_node} base action, but with different occurrence frequencies and degradation intensities. This configuration generates degradation episodes with distinct stochastic characteristics.



\begin{lstlisting}[escapechar=\%,caption={YAML configuration of two independent stocasthic regimes for degrade node event.},label=lst:yaml_tworegimes, basicstyle=\fontsize{8}{10}\selectfont\ttfamily]
infrastructure:
  node:
    actions:
      low_frequent_degrade:                       # first stochastic regime for degrade_node 
        frequency:
          type: exponential
          scale: 30                               # very usual event
        impact:
          composed_of:
            - action_type: degrade_node
              type_object: graph_node
              impact_params:
                target_resolution:
                  mode: self
                target_attributes: ["cpu"]
                p_loss_dist:
                  type: uniform.                  # low degradation level
                  low: 0.10
                  high: 0.25
                distribution_to_restore_node:
                  type: exponential               # short MTTR
                  scale: 3
      high_hardlyever_degrade:                    # second stochastic regime for degrade_node
        frequency:
          type: exponential
          scale: 3000                             # high inter-arrival time
        impact:
          composed_of:
            - action_type: degrade_node
              type_object: graph_node
              impact_params:
                target_resolution:
                  mode: self
                target_attributes: ["cpu"]
                p_loss_dist:
                  type: uniform.                  # high degradation level
                  low: 0.80
                  high: 0.95
                distribution_to_restore_node:
                  type: exponential               # long MTTR
                  scale: 300
\end{lstlisting}


\bmsection{YAML Configuration File Schema for Solver Definition\label{append:solvertriggerpolicy}}

This appendix provides an exhaustive list of examples for the different \textsc{ECCOS} solver trigger policies. Each example illustrates how the corresponding policy is expressed in the YAML solver configuration and which auxiliary parameters are required to control when placement recalculation is executed during the simulation.

\begin{lstlisting}[caption={YAML configuration examples list for all the solver trigger policy alternatives.},label=lst:yaml_tworegimes, basicstyle=\fontsize{8}{10}\selectfont\ttfamily]

trigger_policy:
  type: solve_all

trigger_policy:
  type: solve_none

trigger_policy:
  type: solve_every_t_seconds
  interval_seconds: 30.0

trigger_policy:
  type: solve_every_n_events
  batch_size: 25

trigger_policy:
  type: solve_on_event_types
  critical_events:
    - "disable_node"
    - "disable_edge"
    - "surge_popularity"

trigger_policy:
  type: combined
  interval_seconds: 60.0
  critical_events:
    - "disable_node"
    - "disable_edge"
    - "surge_popularity"

trigger_policy:
  type: solve_random_prob
  probability: 0.10                              # 10\% of probability to execute the solver after each event

trigger_policy:
  type: solve_time_windows
  windows:
    - [0.0, 100.0]
    - [500.0, 600.0]

trigger_policy:
  type: solve_event_index_ranges
  ranges:
    - [1, 50]
    - [200, 250]

trigger_policy:
  type: solve_custom_pattern
  pattern: [true, false, false, false]           # Executes the solver exactly once every 4 events
\end{lstlisting}


%\begin{center}
%\begin{table*}[!ht]%
%\caption{Exhaustive catalog of statistical distribution configurations across simulation domains and their behavioral representations.\label{tab:stochastic_distributions}}
%\begin{tabular*}{\textwidth}{@{\extracolsep\fill}l p{6.0cm} p{7.8cm}@{\extracolsep\fill}}
%\toprule
%&\multicolumn{1}{@{}l}{\textbf{Configuration Target$^{\tnote{\rm a}}$}} & \multicolumn{1}{@{}l}{\textbf{Behavioral Modeling$^{\tnote{\rm b}}$}} \\\cmidrule{2-2}\cmidrule{3-3}
%\textbf{Domain} & \textbf{YAML Attribute Path}  & \textbf{Simulation Representation \& Semantic Meaning}   \\
%\midrule
%\multirow{10}{*}{\textbf{\texttt{infrastructure}}} 
% & \texttt{node.attributes.<attr>.distribution} & Initial computational capacity distribution (CPU, RAM, storage) assigned to physical servers at initialization. \\
% & \texttt{node.attributes.<attr>.distributions.<tier>} & Tier-specific resource capacity distributions (Cloud, Fog, Edge layers). \\
% & \texttt{edge.attributes.<attr>.distribution} & Initial network link capacity distributions (delay, bandwidth, packet loss). \\
% & \texttt{node.actions.<action>.frequency} & Inter-arrival time distribution for stochastic server events (\texttt{disable\_node}, \texttt{degrade\_node}). \\
% & \texttt{node.actions.disable\_node...enable\_node}\tnote{1} & Downtime duration before a disabled server is automatically revived (\texttt{revive\_node}). \\
% & \texttt{node.actions.degrade\_node...p\_loss\_dist}\tnote{2} & Percentage capacity reduction inflicted on a degraded server. \\
% & \texttt{node.actions.degrade\_node...restore\_node}\tnote{3} & Degradation duration before full server capacity is restored. \\
% & \texttt{edge.actions.<action>.frequency} & Inter-arrival time distribution for network link events (\texttt{disable\_edge}, \texttt{congest\_edge}). \\
% & \texttt{edge.actions.disable\_edge...enable\_edge}\tnote{4} & Link downtime duration before connection restoration (\texttt{revive\_edge}). \\
% & \texttt{edge.actions.congest\_edge...p\_loss\_dist}\tnote{5} & Transient packet loss rate or bandwidth reduction during link congestion. \\
% & \texttt{edge.actions.congest\_edge...multiplier\_dist}\tnote{6} & Latency multiplier or bandwidth scaling factor during congestion. \\
% & \texttt{edge.actions.congest\_edge...clear\_edge}\tnote{7} & Congestion duration before normal link parameters are restored. \\
%\midrule
%\multirow{13}{*}{\textbf{\texttt{app}}} 
% & \texttt{services.<attr>.<metric>.distribution} & Initial computational demands (CPU, RAM) or network constraints (delay, bw) assigned to microservices. \\
% & \texttt{popularity.distribution} & Global popularity distribution (e.g., Zipf, Pareto) governing application selection probability. \\
% & \texttt{num\_new\_users} & Distribution governing initial user instances spawned per newly created application. \\
% & \texttt{actions.<action>.frequency} & Inter-arrival time distribution for application events (\texttt{surge}, \texttt{drop}, \texttt{shift}, \texttt{footprint}, \texttt{network}, \texttt{topology}, \texttt{remove}). \\
% & \texttt{actions.surge\_popularity...transient\_prob}\tnote{8} & Probability distribution determining if a popularity surge is transient or permanent. \\
% & \texttt{actions.surge\_popularity...duration\_dist}\tnote{9} & Duration distribution of a transient popularity spike before returning to normal. \\
% & \texttt{actions.drop\_popularity...transient\_prob}\tnote{10} & Probability distribution determining if a popularity drop is transient or permanent. \\
% & \texttt{actions.drop\_popularity...duration\_dist}\tnote{11} & Duration distribution of a transient popularity drop. \\
% & \texttt{actions.geo\_demand\_shift...jump\_prob}\tnote{12} & Probability distribution governing abrupt geographic jumps vs. smooth drift movements. \\
% & \texttt{actions.geo\_demand\_shift...drift\_dist}\tnote{13} & Spatial displacement distribution $(\Delta X, \Delta Y)$ applied during smooth geographic demand drifts. \\
% & \texttt{actions.update\_app\_footprint...multiplier}\tnote{14} & Scaling multiplier distribution applied to microservice computational demands. \\
% & \texttt{actions.update\_app\_footprint...sigma\_...}\tnote{15} & Noise distribution/standard deviation for stochastic resource requirement variations. \\
% & \texttt{actions.update\_app\_network...multiplier}\tnote{16} & Scaling factor distribution applied to bandwidth requirements or delay thresholds. \\
% & \texttt{actions.update\_app\_network...sigma\_net}\tnote{17} & Noise distribution for stochastic network constraint variations. \\
% & \texttt{actions.update\_app\_topology...l\_min / l\_max}\tnote{18} & Distributions defining structural bounds for microservice DAG addition/removal. \\
%\midrule
%\multirow{6}{*}{\textbf{\texttt{user}}} 
% & \texttt{mobility.speed} & Velocity distribution of mobile users traversing the spatial region. \\
% & \texttt{mobility.coverage\_radius} & Wireless transmission radius distribution defining user-to-Edge association thresholds. \\
% & \texttt{request\_ratio} & Initial workload generation rate (request multiplier) assigned to newly created users. \\
% & \texttt{actions.<action>.frequency} & Inter-arrival time distribution for user stochastic events (\texttt{move}, \texttt{change\_ratio}, \texttt{suspend}, \texttt{remove}). \\
% & \texttt{actions.change\_request\_ratio...multiplier}\tnote{19} & Scaling factor distribution applied to dynamically modify a user's request ratio. \\
% & \texttt{actions.suspend\_user...resume\_user}\tnote{20} & Session suspension duration before a user automatically resumes traffic generation. \\
%\midrule
%\multirow{2}{*}{\textbf{\texttt{global}}} 
% & \texttt{global\_spawner.actions.new\_user.frequency} & Inter-arrival time distribution for dynamic \textit{ex novo} user spawning. \\
% & \texttt{global\_spawner.actions.new\_app.frequency} & Inter-arrival time distribution for dynamic \textit{ex novo} application and user spawning. \\
%\bottomrule
%\end{tabular*}
%\begin{tablenotes}%%[341pt]
%\item[$^{\rm a}$] Exact YAML hierarchical attribute path where a distribution dictionary (\texttt{type}, \texttt{low/high}, \texttt{mean/sigma}, etc.) is evaluated.
%\item[$^{\rm b}$] Semantic behavioral effect in the simulation engine.
%\item[1] Full path: \texttt{node.actions.disable\_node.impact.distribution\_to\_enable\_node}.
%\item[2] Full path: \texttt{node.actions.degrade\_node.impact.p\_loss\_dist}.
%\item[3] Full path: \texttt{node.actions.degrade\_node.impact.distribution\_to\_restore\_node}.
%\item[4] Full path: \texttt{edge.actions.disable\_edge.impact.distribution\_to\_enable\_edge}.
%\item[5] Full path: \texttt{edge.actions.congest\_edge.impact.p\_loss\_dist}.
%\item[6] Full path: \texttt{edge.actions.congest\_edge.impact.multiplier\_dist}.
%\item[7] Full path: \texttt{edge.actions.congest\_edge.impact.distribution\_to\_clear\_edge}.
%\item[8] Full path: \texttt{actions.surge\_popularity.impact.transient\_prob}.
%\item[9] Full path: \texttt{actions.surge\_popularity.impact.duration\_dist}.
%\item[10] Full path: \texttt{actions.drop\_popularity.impact.transient\_prob}.
%\item[11] Full path: \texttt{actions.drop\_popularity.impact.duration\_dist}.
%\item[12] Full path: \texttt{actions.geo\_demand\_shift.impact.jump\_prob}.
%\item[13] Full path: \texttt{actions.geo\_demand\_shift.impact.drift\_dist}.
%\item[14] Full path: \texttt{actions.update\_app\_footprint.impact.multiplier}.
%\item[15] Full path: \texttt{actions.update\_app\_footprint.impact.sigma\_footprint}.
%\item[16] Full path: \texttt{actions.update\_app\_network.impact.multiplier}.
%\item[17] Full path: \texttt{actions.update\_app\_network.impact.sigma\_net}.
%\item[18] Full path: \texttt{actions.update\_app\_topology.impact.l\_min} and \texttt{l\_max\_global}.
%\item[19] Full path: \texttt{actions.change\_request\_ratio.impact.multiplier}.
%\item[20] Full path: \texttt{actions.suspend\_user.impact.distribution\_to\_resume\_user}.
%\item {\it Source:} Compiled from the simulation engine distribution parser (\texttt{parse\_distribution}) and YAML schema definitions.
%\end{tablenotes}
%\end{table*}
%\end{center}




\end{document}
